\chapter{Groups I}
\lecture{1}{Rotman §2.2--§2.4}{Permutations, Groups, Lagrange's Theorem}

\section{Permutations}

In high school mathematics the words \emph{permutation} and \emph{arrangement} are used
interchangeably. We draw a distinction between them.

\begin{definition}[Permutation, arrangement]
	A \vocab{permutation} of a set \(X\) is a bijection \(\alpha\colon X \to X\). If \(X\)
	is finite with \(\abs{X} = n\), then an \vocab{arrangement} of \(X\) is a list
	\(x_1, x_2, \dots, x_n\) of all the elements of \(X\) with no repetitions.
\end{definition}

Given an arrangement \(x_1, \dots, x_n\), define \(f\colon \{1, \dots, n\} \to X\) by
\(f(i) = x_i\); thus the list displays the values of \(f\). That there are no repetitions
says \(f\) is injective, and that every \(x \in X\) occurs on the list says \(f\) is
surjective. So an arrangement of \(X\) is the same thing as a bijection
\(f\colon \{1, \dots, n\} \to X\).

\begin{eg}[Arrangements of a three-element set]
	There are six arrangements of \(X = \{a,b,c\}\):
	\[
		abc; \quad acb; \quad bac; \quad bca; \quad cab; \quad cba.
	\]
	In general there are exactly \(n!\) arrangements of an \(n\)-element set.
\end{eg}

If \(X = \{1, 2, \dots, n\}\), a permutation \(\alpha\colon X \to X\) gives the list
\(\alpha(1) = i_1, \dots, \alpha(n) = i_n\), which we record in \vocab{two-rowed notation}:
\[
	\alpha =
	\begin{pmatrix}
		1         & 2         & \cdots & j            & \cdots & n         \\
		\alpha(1) & \alpha(2) & \cdots & \alpha(j)    & \cdots & \alpha(n)
	\end{pmatrix}.
\]

\begin{intuition}
	Informally, arrangements (lists) and permutations (bijections) are two ways of
	describing the same thing. The advantage of viewing a permutation as a \emph{bijection}
	is that bijections can be \emph{composed}, and the composite of two permutations is
	again a permutation. Lists can only be counted.
\end{intuition}

\begin{definition}[Symmetric group]
	The family of all permutations of a set \(X\), denoted by \(S_X\), is called the
	\vocab{symmetric group} on \(X\). When \(X = \{1, 2, \dots, n\}\) we write \(S_n\) and
	call it the \vocab{symmetric group on \(n\) letters}.
\end{definition}

\begin{eg}[Composition in \(S_3\) is not commutative]
	If
	\[
		\alpha = \begin{pmatrix} 1 & 2 & 3 \\ 2 & 3 & 1 \end{pmatrix}
		\quad\text{and}\quad
		\beta  = \begin{pmatrix} 1 & 2 & 3 \\ 2 & 1 & 3 \end{pmatrix},
	\]
	then
	\[
		\alpha \circ \beta = \begin{pmatrix} 1 & 2 & 3 \\ 3 & 2 & 1 \end{pmatrix}
		\quad\text{and}\quad
		\beta \circ \alpha = \begin{pmatrix} 1 & 2 & 3 \\ 1 & 3 & 2 \end{pmatrix},
	\]
	so \(\alpha \circ \beta \neq \beta \circ \alpha\): for instance
	\(\alpha\circ\beta\colon 1 \mapsto \alpha(\beta(1)) = \alpha(2) = 3\), whereas
	\(\beta\circ\alpha\colon 1 \mapsto 2 \mapsto 1\).

	Some permutations do commute, however; the reader may check that
	\(\begin{psmallmatrix} 1&2&3&4 \\ 2&1&3&4\end{psmallmatrix}\) and
	\(\begin{psmallmatrix} 1&2&3&4 \\ 1&2&4&3\end{psmallmatrix}\) commute.
\end{eg}

\begin{note}[Order of multiplication]
	Some authors put ``functions on the right'', writing \((i)\alpha\) instead of
	\(\alpha(i)\). Their \(\alpha \circ \beta\) is our \(\beta \circ \alpha\): the
	notational switch causes a switch in the order of multiplication.
\end{note}

\begin{proposition}[Cancellation in \(S_X\)]\label{prop:cancel-SX}
	If \(\gamma, \alpha, \beta \in S_X\) and \(\gamma\alpha = \gamma\beta\), then
	\(\alpha = \beta\). Similarly \(\alpha\gamma = \beta\gamma\) implies \(\alpha = \beta\).
\end{proposition}

\begin{proof}
	Compose with \(\gamma^{-1}\) and use associativity of composition:
	\[
		\alpha
		= 1_X \circ \alpha
		= (\gamma^{-1}\circ\gamma) \circ \alpha
		= \gamma^{-1} \circ (\gamma \circ \alpha)
		= \gamma^{-1} \circ (\gamma \circ \beta)
		= (\gamma^{-1}\circ\gamma) \circ \beta
		= 1_X \circ \beta = \beta.
	\]
	A similar argument works on the other side.
\end{proof}

Aside from being cumbersome, the two-rowed notation hides the answers to elementary
questions: is the square of a permutation the identity? what is the smallest \(m > 0\)
with \(\alpha^m = 1\)? can one factor a permutation into simpler pieces? The special
permutations below remedy this defect.

\begin{notation}
	From now on we write \(\beta\alpha\) instead of \(\beta \circ \alpha\), and \((1)\)
	instead of \(1_X\).
\end{notation}

\subsection{Cycles}

\begin{definition}[Fixes, moves]
	If \(\alpha \in S_n\) and \(i \in \{1, \dots, n\}\), then \(\alpha\) \vocab{fixes} \(i\)
	if \(\alpha(i) = i\), and \(\alpha\) \vocab{moves} \(i\) if \(\alpha(i) \neq i\).
\end{definition}

\begin{definition}[Cycle]
	Let \(i_1, i_2, \dots, i_r\) be distinct integers in \(\{1, \dots, n\}\). If
	\(\alpha \in S_n\) fixes the other integers (if any) and
	\[
		\alpha(i_1) = i_2, \quad \alpha(i_2) = i_3, \quad \dots, \quad
		\alpha(i_{r-1}) = i_r, \quad \alpha(i_r) = i_1,
	\]
	then \(\alpha\) is called an \vocab{\(r\)-cycle}, or a \vocab{cycle of length \(r\)},
	and we write
	\[
		\alpha = (i_1\ i_2\ \dots\ i_r).
	\]
	A \(2\)-cycle is called a \vocab{transposition}. A \(1\)-cycle is the identity, so all
	\(1\)-cycles are equal: \((i) = (1)\) for every \(i\).
\end{definition}

\begin{eg}[Cycle notation]
	The two-rowed notation does not help us recognize that
	\(\alpha = \begin{psmallmatrix} 1&2&3&4&5 \\ 4&3&1&5&2 \end{psmallmatrix}\) is a
	\(5\)-cycle, but in fact \(\alpha(1) = 4\), \(\alpha(4) = 5\), \(\alpha(5) = 2\),
	\(\alpha(2) = 3\), \(\alpha(3) = 1\); that is, \(\alpha = (1\ 4\ 5\ 2\ 3)\). Likewise
	\begin{gather*}
		\begin{pmatrix} 1&2&3&4 \\ 2&3&4&1 \end{pmatrix} = (1\ 2\ 3\ 4),
		\qquad
		\begin{pmatrix} 1&2&3&4&5 \\ 5&1&4&2&3 \end{pmatrix} = (1\ 5\ 3\ 4\ 2), \\
		\begin{pmatrix} 1&2&3&4&5 \\ 2&3&1&4&5 \end{pmatrix} = (1\ 2\ 3).
	\end{gather*}
	On the other hand \(\beta = \begin{psmallmatrix} 1&2&3&4 \\ 2&1&4&3\end{psmallmatrix}\)
	is \emph{not} a cycle; in fact \(\beta = (1\ 2)(3\ 4)\).
\end{eg}

The term \emph{cycle} comes from the Greek word for circle: one pictures the cycle
\((i_1\ i_2\ \dots\ i_r)\) as a clockwise rotation of a circle
(\autoref{fig:cycle-rotation}). Any \(i_j\) may be taken as the starting point, so an
\(r\)-cycle has \(r\) different cycle notations:
\[
	(i_1\ i_2\ \dots\ i_r) = (i_2\ i_3\ \dots\ i_r\ i_1) = \cdots = (i_r\ i_1\ i_2\ \dots\ i_{r-1}).
\]

\begin{figure}[H]
	\centering
	\begin{tikzpicture}[>=stealth]
		\def\rad{1.9}
		\draw[gray!40, dashed] (0,0) circle (\rad);
		\fill (90:\rad)  circle (1.8pt) node[above=3pt]      {\(i_1\)};
		\fill (20:\rad)  circle (1.8pt) node[right=3pt]      {\(i_2\)};
		\fill (-50:\rad) circle (1.8pt) node[below right=1pt] {\(i_3\)};
		\fill (160:\rad) circle (1.8pt) node[left=3pt]       {\(i_r\)};
		\foreach \a in {272, 245, 218, 191} {
			\fill[black!75] (\a:\rad) circle (1.3pt);
		}
		\draw[->, thick, MidnightBlue] (82:\rad) arc (82:28:\rad);
		\draw[->, thick, MidnightBlue] (12:\rad) arc (12:-42:\rad);
		\draw[->, thick, MidnightBlue] (152:\rad) arc (152:98:\rad);
	\end{tikzpicture}
	\caption{A cycle as a (clockwise) rotation.}
	\label{fig:cycle-rotation}
\end{figure}

\begin{eg}[The factorization algorithm]\label{eg:factor-algorithm}
	Take
	\[
		\alpha = \begin{pmatrix}
			1&2&3&4&5&6&7&8&9 \\
			6&4&7&2&5&1&8&9&3
		\end{pmatrix}.
	\]
	Begin by writing ``\((1\)''. Now \(\alpha\colon 1 \mapsto 6\), so write ``\((1\ 6\)''.
	Next \(\alpha\colon 6 \mapsto 1\), so the parentheses close: \(\alpha\) begins
	``\((1\ 6)\)''. The smallest number not yet used is \(2\), and \(2 \mapsto 4\),
	\(4 \mapsto 2\), giving ``\((1\ 6)(2\ 4)\)''. The smallest remaining number is \(3\),
	and \(3 \mapsto 7 \mapsto 8 \mapsto 9 \mapsto 3\), giving the \(4\)-cycle
	\((3\ 7\ 8\ 9)\). Finally \(\alpha(5) = 5\), and so
	\[
		\alpha = (1\ 6)(2\ 4)(3\ 7\ 8\ 9)(5).
	\]
	To verify a claim like this one evaluates both sides. For instance, writing
	\(\beta = (1\ 6)\), \(\gamma = (2\ 4)\), \(\delta = (3\ 7\ 8\ 9)\),
	\[
		\beta\gamma\delta(1) = \beta(\gamma(\delta(1))) = \beta(\gamma(1)) = \beta(1) = 6 = \alpha(1),
	\]
	since \(\delta\) and \(\gamma\) fix \(1\). (\autoref{prop:disjoint-cycles} gives a more
	satisfactory proof that the algorithm works.)
\end{eg}

\begin{eg}[Multiplying with the algorithm]
	Factorizations into cycles are convenient for \emph{multiplying} permutations: one
	simply displays the partial outputs. In \(S_5\) let
	\[
		\sigma = (1\ 2)(1\ 3\ 4\ 2\ 5)(2\ 5\ 1\ 3).
	\]
	Then \(\sigma\colon 1 \mapsto 3 \mapsto 4 \mapsto 4\), so \(\sigma\) begins
	``\((1\ 4\)''. Next \(4 \mapsto 4 \mapsto 2 \mapsto 1\), so \(\sigma\) begins
	``\((1\ 4)\)''. Then \(2 \mapsto 5 \mapsto 1 \mapsto 2\), so \(\sigma\) fixes \(2\);
	and \(3 \mapsto 2 \mapsto 5 \mapsto 5\), \(5 \mapsto 1 \mapsto 3 \mapsto 3\).
	Therefore \(\sigma = (1\ 4)(2)(3\ 5)\).
\end{eg}

\begin{definition}[Disjoint permutations]
	Two permutations \(\alpha, \beta \in S_n\) are \vocab{disjoint} if every \(i\) moved by
	one is fixed by the other: if \(\alpha(i) \neq i\) then \(\beta(i) = i\), and if
	\(\beta(j) \neq j\) then \(\alpha(j) = j\). A family \(\beta_1, \dots, \beta_t\) is
	disjoint if each pair of them is disjoint.
\end{definition}

\begin{observation}
	If \(\alpha = (i_1\ \dots\ i_r)\) and \(\beta = (j_1\ \dots\ j_s)\), then any
	\(k \in \{i_1, \dots, i_r\} \cap \{j_1, \dots, j_s\}\) is moved by both. Hence two
	cycles are disjoint if and only if the underlying sets
	\(\{i_1, \dots, i_r\}\) and \(\{j_1, \dots, j_s\}\) are disjoint sets. When
	\(\alpha\) and \(\beta\) are disjoint there are exactly three possibilities for a
	number \(i\): it is moved by \(\alpha\), it is moved by \(\beta\), or it is fixed by
	both.
\end{observation}

\begin{lemma}[Disjoint permutations commute]\label{lem:disjoint-commute}
	Disjoint permutations \(\alpha, \beta \in S_n\) commute.
\end{lemma}

\begin{proof}
	It suffices to prove \(\alpha\beta(i) = \beta\alpha(i)\) for \(1 \leq i \leq n\).
	If \(\beta\) moves \(i\), say \(\beta(i) = j \neq i\), then \(\beta\) also moves \(j\)
	(otherwise \(\beta(j) = j\) and \(\beta(i) = j\) contradict injectivity of \(\beta\));
	since \(\alpha\) and \(\beta\) are disjoint, \(\alpha(i) = i\) and \(\alpha(j) = j\),
	so \(\beta\alpha(i) = \beta(i) = j = \alpha(j) = \alpha\beta(i)\). A similar argument
	applies if \(\alpha\) moves \(i\). The last possibility is that neither moves \(i\),
	and then \(\alpha\beta(i) = i = \beta\alpha(i)\).
\end{proof}

\begin{remark}
	The converse fails: permutations that are \emph{not} disjoint may still commute. The
	reader may check that \((1\ 2\ 3)(4\ 5)\) and \((1\ 3\ 2)(6\ 7)\) commute; even simpler,
	\(\alpha\alpha^2 = \alpha^2\alpha\).
\end{remark}

\begin{lemma}\label{lem:orbit}
	Let \(X = \{1, \dots, n\}\), let \(\alpha \in S_X = S_n\), and given \(i_1 \in X\)
	define \(i_{j+1} = \alpha(i_j)\) for \(j \geq 1\). Write \(Y = \{i_j : j \geq 1\}\) and
	let \(Y'\) be the complement of \(Y\) in \(X\).
	\begin{enumerate}[(i)]
		\item If \(\alpha\) moves \(i_1\), then there is \(r > 1\) with
		      \(i_1, \dots, i_r\) all distinct and \(i_{r+1} = \alpha(i_r) = i_1\).
		\item \(\alpha(Y) = Y\) and \(\alpha(Y') = Y'\).
	\end{enumerate}
\end{lemma}

\begin{proof}
	(i) Since \(X\) is finite, there is a smallest \(r > 1\) with \(i_1, \dots, i_r\) all
	distinct but \(i_{r+1} = \alpha(i_r) \in \{i_1, \dots, i_r\}\); say
	\(\alpha(i_r) = i_j\) with \(1 \leq j \leq r\). If \(j > 1\), then
	\(\alpha(i_r) = i_j = \alpha(i_{j-1})\), and injectivity of \(\alpha\) gives
	\(i_r = i_{j-1}\), contradicting distinctness. Therefore \(\alpha(i_r) = i_1\).

	(ii) Clearly \(\alpha(Y) \subseteq Y\), since \(i_j \in Y\) implies
	\(\alpha(i_j) = i_{j+1} \in Y\). If \(k \in Y'\) and \(\alpha(k) \in Y\), then
	\(\alpha(k) = i_j = \alpha(i_{j-1})\) for some \(j\) (by part (i) this holds even for
	\(i_j = i_1\)); injectivity gives \(k = i_{j-1} \in Y\), contradicting
	\(Y \cap Y' = \varnothing\). Hence \(\alpha(Y') \subseteq Y'\).

	These inclusions are equalities: \(\alpha(X) = \alpha(Y) \cup \alpha(Y')\) is a disjoint
	union because \(\alpha\) is injective, and \(\abs{\alpha(Y)} \leq \abs{Y}\),
	\(\abs{\alpha(Y')} \leq \abs{Y'}\). If either inequality were strict, then
	\(\abs{\alpha(X)} < \abs{X}\), contradicting surjectivity of \(\alpha\).
\end{proof}

\begin{proposition}[Factorization into disjoint cycles]\label{prop:disjoint-cycles}
	Every permutation \(\alpha \in S_n\) is either a cycle or a product of disjoint cycles.
\end{proposition}

\begin{proof}
	Induct on the number \(k \geq 0\) of points moved by \(\alpha\). The base step
	\(k = 0\) is true, for then \(\alpha\) is the identity, which is a \(1\)-cycle.

	If \(k > 0\), let \(i_1\) be a point moved by \(\alpha\). As in \autoref{lem:orbit},
	define \(Y = \{i_1, \dots, i_r\}\) with the \(i_j\) distinct, \(\alpha(i_j) = i_{j+1}\)
	for \(j < r\), and \(\alpha(i_r) = i_1\). Let \(\sigma = (i_1\ i_2\ \dots\ i_r)\), so
	that \(\sigma\) fixes each point of \(Y'\). If \(r = n\), then \(\alpha = \sigma\).
	If \(r < n\), then \(\alpha(Y') = Y'\) by the lemma. Put \(\alpha' = \alpha\sigma^{-1}\);
	we claim \(\alpha'\) and \(\sigma\) are disjoint. If \(\sigma\) moves \(i\), then
	\(i = i_j \in Y\), and
	\(\alpha'(i_j) = \alpha\sigma^{-1}(i_j) = \alpha(i_{j-1}) = i_j\), so \(\alpha'\) fixes
	\(i_j\). Conversely, if \(\alpha'\) moves \(k\), then \(k \notin Y\), so \(k \in Y'\);
	but \(\sigma\) fixes every element of \(Y'\). Hence \(\alpha = \alpha'\sigma\) is a
	factorization into disjoint permutations. The number of points moved by \(\alpha'\) is
	\(k - r < k\), so the inductive hypothesis gives \(\alpha' = \beta_1 \cdots \beta_t\)
	with \(\beta_1, \dots, \beta_t\) disjoint cycles. Therefore
	\(\alpha = \beta_1 \cdots \beta_t \sigma\) is a product of disjoint cycles.
\end{proof}

\begin{definition}[Complete factorization]
	A \vocab{complete factorization} of a permutation \(\alpha\) is a factorization of
	\(\alpha\) into disjoint cycles that contains one \(1\)-cycle \((i)\) for every \(i\)
	fixed by \(\alpha\).
\end{definition}

The algorithm of \autoref{eg:factor-algorithm} always yields a complete factorization. In
a complete factorization \(\alpha = \beta_1\cdots\beta_t\), every symbol \(i\) between
\(1\) and \(n\) occurs in exactly one \(\beta_\lambda\).

\begin{eg}
	If \(\alpha = \begin{psmallmatrix} 1&2&3&4&5 \\ 1&3&4&2&5 \end{psmallmatrix}\), the
	algorithm gives the complete factorization \(\alpha = (1)(2\ 3\ 4)(5)\). By contrast the
	factorizations
	\[
		\alpha = (2\ 3\ 4) = (1)(2\ 3\ 4) = (2\ 3\ 4)(5)
	\]
	are \emph{not} complete.
\end{eg}

Write \(\beta = (i_0\ i_1\ \dots\ i_{r-1})\). Then \(i_1 = \beta(i_0)\),
\(i_2 = \beta(i_1) = \beta^2(i_0)\), and in general
\begin{equation}\label{eq:cycle-power}
	i_k = \beta^k(i_0) \qquad (k \leq r-1).
\end{equation}
Since \(\beta(i_{r-1}) = i_0\), equation~\eqref{eq:cycle-power} holds for all \(k\) if the
subscripts \(j\) in \(i_j\) are read \(\bmod\ r\).

\begin{lemma}[Powers of a cycle]\label{lem:cycle-order}
	Let \(\beta = (i_0\ i_1\ \dots\ i_{r-1})\) be an \(r\)-cycle. Then \(\beta^r = (1)\),
	and \(r\) is the smallest positive integer \(k\) with \(\beta^k = (1)\).
\end{lemma}

\begin{proof}
	Reading subscripts \(\bmod\ r\), equation~\eqref{eq:cycle-power} gives
	\(\beta^k(i_j) = i_{j+k}\) for all \(j\) and all \(k \geq 1\). Taking \(k = r\) yields
	\(\beta^r(i_j) = i_j\) for every \(j\); since \(\beta\) fixes every other symbol, so
	does \(\beta^r\), and hence \(\beta^r = (1)\). If \(0 < k < r\), then
	\(\beta^k(i_0) = i_k \neq i_0\), because \(i_0, \dots, i_{r-1}\) are distinct;
	therefore \(\beta^k \neq (1)\).
\end{proof}

\begin{lemma}\label{lem:cycle-powers}
	\begin{enumerate}[(i)]
		\item Let \(\alpha = \beta\delta\) be a factorization into disjoint permutations. If
		      \(\beta\) moves \(i\), then \(\alpha^k(i) = \beta^k(i)\) for all \(k \geq 1\).
		\item If \(\beta\) and \(\gamma\) are cycles both of which move \(i = i_0\), and if
		      \(\beta^k(i) = \gamma^k(i)\) for all \(k \geq 1\), then \(\beta = \gamma\).
	\end{enumerate}
\end{lemma}

\begin{remark}
	The hypothesis in (ii) does \emph{not} assume that \(\beta\) and \(\gamma\) have the
	same length; that is part of the conclusion.
\end{remark}

\begin{proof}
	(i) Since \(\beta\) moves \(i\), disjointness implies \(\delta\) fixes \(i\), and hence
	so does every power of \(\delta\). Now \(\beta\) and \(\delta\) commute by
	\autoref{lem:disjoint-commute}, so
	\((\beta\delta)^k(i) = \beta^k(\delta^k(i)) = \beta^k(i)\).

	(ii) By \eqref{eq:cycle-power}, if \(\beta = (i_0\ i_1\ \dots\ i_{r-1})\), then
	\(i_k = \beta^k(i_0)\), and if \(\gamma = (i_0\ j_1\ \dots\ j_{s-1})\), then
	\(j_k = \gamma^k(i_0)\). We may assume \(r \leq s\), so that
	\(i_1 = j_1, \dots, i_{r-1} = j_{r-1}\). Since
	\(j_r = \gamma^r(i_0) = \beta^r(i_0) = i_0\), it follows that \(s = r\) and
	\(j_k = i_k\) for all \(k\). Therefore \(\beta = \gamma\).
\end{proof}

\begin{theorem}[Uniqueness of the complete factorization]\label{thm:unique-factorization}
	Let \(\alpha \in S_n\) and let \(\alpha = \beta_1 \cdots \beta_t\) be a complete
	factorization into disjoint cycles. This factorization is unique except for the order in
	which the cycles occur.
\end{theorem}

\begin{intuition}
	This is the analog, for permutations, of the fundamental theorem of arithmetic: cycles
	play the role of primes.
\end{intuition}

\begin{proof}
	Let \(\alpha = \gamma_1 \cdots \gamma_s\) be a second complete factorization into
	disjoint cycles. Every complete factorization has exactly one \(1\)-cycle for each \(i\)
	fixed by \(\alpha\), so it suffices to prove, by induction on \(\ell\), the larger of
	\(t\) and \(s\), that the cycles of length \(> 1\) are uniquely determined by \(\alpha\).

	The base step \(\ell = 1\) is the hypothesis \(\beta_1 = \alpha = \gamma_1\).

	For the inductive step, if \(\beta_t\) moves \(i = i_0\), then
	\(\beta_t^{\,k}(i_0) = \alpha^k(i_0)\) for all \(k \geq 1\), by
	\autoref{lem:cycle-powers}(i). Some \(\gamma_j\) must also move \(i_0\); since disjoint
	cycles commute, we may re-index so that \(\gamma_s\) moves \(i_0\), and then
	\(\gamma_s^{\,k}(i_0) = \alpha^k(i_0)\) for all \(k\). Now
	\autoref{lem:cycle-powers}(ii) gives \(\beta_t = \gamma_s\), so that cancellation
	(\autoref{prop:cancel-SX}) yields
	\[
		\beta_1 \cdots \beta_{t-1} = \gamma_1 \cdots \gamma_{s-1}.
	\]
	By the inductive
	hypothesis \(s = t\) and the \(\gamma\)'s may be re-indexed so that
	\(\gamma_\lambda = \beta_\lambda\) for all \(\lambda\).
\end{proof}

\begin{proposition}[Inverses]\label{prop:inverse-cycle}
	\begin{enumerate}[(i)]
		\item \((i_1\ i_2\ \dots\ i_r)^{-1} = (i_r\ i_{r-1}\ \dots\ i_1)\).
		\item If \(\gamma \in S_n\) and \(\gamma = \beta_1 \cdots \beta_k\), then
		      \(\gamma^{-1} = \beta_k^{-1} \cdots \beta_1^{-1}\) (note that the order of the
		      factors is reversed).
	\end{enumerate}
\end{proposition}

\begin{intuition}
	In the picture of a cycle as a clockwise rotation of a circle
	(\autoref{fig:cycle-rotation}), the inverse is just the counterclockwise rotation.
\end{intuition}

\begin{proof}
	(i) The composite \((i_1\ i_2\ \dots\ i_r)(i_r\ i_{r-1}\ \dots\ i_1)\) fixes every
	integer between \(1\) and \(n\) other than \(i_1, \dots, i_r\); it sends
	\(i_1 \mapsto i_r \mapsto i_1\), and for \(j \geq 2\) it sends
	\(i_j \mapsto i_{j-1} \mapsto i_j\). Thus every integer is fixed, so the composite is
	\((1)\). A similar argument treats the other order.

	(ii) Induct on \(k \geq 2\). For \(k = 2\),
	\[
		(\beta_1\beta_2)(\beta_2^{-1}\beta_1^{-1})
		= \beta_1(\beta_2\beta_2^{-1})\beta_1^{-1}
		= \beta_1\beta_1^{-1} = (1),
	\]
	and similarly \((\beta_2^{-1}\beta_1^{-1})(\beta_1\beta_2) = (1)\). For the inductive
	step, let \(\delta = \beta_1\cdots\beta_k\); then
	\(\beta_1 \cdots \beta_k\beta_{k+1} = \delta\beta_{k+1}\), and so
	\[
		(\beta_1 \cdots \beta_{k+1})^{-1}
		= (\delta\beta_{k+1})^{-1}
		= \beta_{k+1}^{-1}\delta^{-1}
		= \beta_{k+1}^{-1}(\beta_1\cdots\beta_k)^{-1}
		= \beta_{k+1}^{-1}\beta_k^{-1}\cdots\beta_1^{-1}. \qedhere
	\]
\end{proof}

Thus \((1\ 2\ 3\ 4)^{-1} = (4\ 3\ 2\ 1) = (1\ 4\ 3\ 2)\), and
\((1\ 2)^{-1} = (2\ 1) = (1\ 2)\): every transposition is its own inverse.

\begin{eg}
	\autoref{prop:inverse-cycle} applies in particular when the factors are disjoint
	cycles, in which case reversing the order is unnecessary because disjoint cycles commute.
	Thus if \(\alpha = \begin{psmallmatrix} 1&2&3&4&5&6&7&8&9 \\ 6&4&7&2&5&1&8&9&3
	\end{psmallmatrix} = (1\ 6)(2\ 4)(3\ 7\ 8\ 9)(5)\), then
	\[
		\alpha^{-1} = (5)(9\ 8\ 7\ 3)(4\ 2)(6\ 1) = (1\ 6)(2\ 4)(3\ 9\ 8\ 7).
	\]
\end{eg}

\subsection{Cycle Structure and Conjugation}

\begin{definition}[Cycle structure]
	Two permutations \(\alpha, \beta \in S_n\) have the same \vocab{cycle structure} if
	their complete factorizations have the same number of \(r\)-cycles for every
	\(r \geq 1\).
\end{definition}

There are \(\tfrac{1}{r}\bigl[n(n-1)\cdots(n-r+1)\bigr]\) \(r\)-cycles in \(S_n\)
(Rotman, Exercise 2.21). This formula counts permutations of a given cycle
structure, provided one is careful when several cycles have the same length. For example,
the number of permutations in \(S_4\) of cycle structure \((a\ b)(c\ d)\) is
\[
	\tfrac{1}{2}\Bigl[\tfrac{1}{2}(4\times 3)\Bigr]\cdot\Bigl[\tfrac{1}{2}(2 \times 1)\Bigr] = 3,
\]
the extra factor \(\tfrac{1}{2}\) occurring so that we do not count
\((a\ b)(c\ d) = (c\ d)(a\ b)\) twice. Similarly, the number of permutations in \(S_n\) of
the form \((a\ b)(c\ d)(e\ f)\) is
\[
	\frac{1}{3!\,2^3}\bigl[n(n-1)(n-2)(n-3)(n-4)(n-5)\bigr].
\]

\begin{table}[H]
	\centering
	\begin{minipage}{0.42\linewidth}
		\centering
		\begin{tabular}{lr}
			\toprule
			Cycle structure & Number \\
			\midrule
			\((1)\)          & 1  \\
			\((1\ 2)\)       & 6  \\
			\((1\ 2\ 3)\)    & 8  \\
			\((1\ 2\ 3\ 4)\) & 6  \\
			\((1\ 2)(3\ 4)\) & 3  \\
			\midrule
			                 & 24 \\
			\bottomrule
		\end{tabular}
		\subcaption{Permutations in \(S_4\).}
	\end{minipage}
	\hfill
	\begin{minipage}{0.52\linewidth}
		\centering
		\begin{tabular}{lr}
			\toprule
			Cycle structure     & Number \\
			\midrule
			\((1)\)               & 1   \\
			\((1\ 2)\)            & 10  \\
			\((1\ 2\ 3)\)         & 20  \\
			\((1\ 2\ 3\ 4)\)      & 30  \\
			\((1\ 2\ 3\ 4\ 5)\)   & 24  \\
			\((1\ 2)(3\ 4\ 5)\)   & 20  \\
			\((1\ 2)(3\ 4)\)      & 15  \\
			\midrule
			                      & 120 \\
			\bottomrule
		\end{tabular}
		\subcaption{Permutations in \(S_5\).}
	\end{minipage}
	\caption{Counting permutations by cycle structure.}
	\label{tab:cycle-counts}
\end{table}

\begin{lemma}\label{lem:conj-action}
	Let \(\alpha, \gamma \in S_n\). For all \(i\), if \(\gamma\colon i \mapsto j\), then
	\(\alpha\gamma\alpha^{-1}\colon \alpha(i) \mapsto \alpha(j)\).
\end{lemma}

\begin{proof}
	\(\alpha\gamma\alpha^{-1}(\alpha(i)) = \alpha\gamma(i) = \alpha(j)\).
\end{proof}

\begin{proposition}\label{prop:conj-same-structure}
	If \(\gamma, \alpha \in S_n\), then \(\alpha\gamma\alpha^{-1}\) has the same cycle
	structure as \(\gamma\). In more detail, if the complete factorization of \(\gamma\) is
	\[
		\gamma = \beta_1\beta_2 \cdots (i\ j\ \dots) \cdots \beta_t,
	\]
	then \(\alpha\gamma\alpha^{-1}\) is the permutation obtained from \(\gamma\) by applying
	\(\alpha\) to every symbol occurring in the cycles of \(\gamma\).
\end{proposition}

\begin{remark}
	For example, if \(\gamma = (1\ 3)(2\ 4\ 7)(5)(6)\) and
	\(\alpha = (2\ 5\ 6)(1\ 4\ 3)\), then
	\[
		\alpha\gamma\alpha^{-1}
		= (\alpha 1\ \ \alpha 3)(\alpha 2\ \ \alpha 4\ \ \alpha 7)(\alpha 5)(\alpha 6)
		= (4\ 1)(5\ 3\ 7)(6)(2).
	\]
\end{remark}

\begin{proof}
	Let \(\sigma\) denote the permutation obtained from \(\gamma\) by applying \(\alpha\) to
	all the symbols. If \(\gamma\) fixes \(i\), then \autoref{lem:conj-action} shows that
	\(\alpha\gamma\alpha^{-1}\) fixes \(\alpha(i)\), as does \(\sigma\). Assume that
	\(\gamma\) moves \(i\), say \(\gamma(i) = j\), so that one of the cycles of \(\gamma\)
	is \((i\ j\ \dots)\); by definition, one of the cycles of \(\sigma\) is
	\((\alpha(i)\ \alpha(j)\ \dots)\), that is,
	\(\sigma\colon \alpha(i) \mapsto \alpha(j)\). But \autoref{lem:conj-action} says
	\(\alpha\gamma\alpha^{-1}\colon \alpha(i) \mapsto \alpha(j)\), so \(\sigma\) and
	\(\alpha\gamma\alpha^{-1}\) agree on every number of the form \(\alpha(i)\). Since
	\(\alpha\) is surjective, every \(k \in X\) has this form, and so
	\(\sigma = \alpha\gamma\alpha^{-1}\).
\end{proof}

\begin{proposition}\label{prop:conjugacy-criterion}
	If \(\gamma, \gamma' \in S_n\), then \(\gamma\) and \(\gamma'\) have the same cycle
	structure if and only if there exists \(\alpha \in S_n\) with
	\(\gamma' = \alpha\gamma\alpha^{-1}\).
\end{proposition}

\begin{proof}
	Sufficiency is \autoref{prop:conj-same-structure}.

	Conversely, assume \(\gamma\) and \(\gamma'\) have the same cycle structure; that is,
	\(\gamma = \beta_1\cdots\beta_t\) and \(\gamma' = \sigma_1\cdots\sigma_t\) are complete
	factorizations with \(\beta_\lambda\) and \(\sigma_\lambda\) of the same length for all
	\(\lambda \leq t\). Write
	\(\beta_\lambda = (i_1^\lambda, \dots, i_{r(\lambda)}^\lambda)\) and
	\(\sigma_\lambda = (j_1^\lambda, \dots, j_{r(\lambda)}^\lambda)\), and define
	\[
		\alpha(i_1^\lambda) = j_1^\lambda, \quad
		\alpha(i_2^\lambda) = j_2^\lambda, \quad \dots, \quad
		\alpha(i_{r(\lambda)}^\lambda) = j_{r(\lambda)}^\lambda
	\]
	for all \(\lambda\). Since \(\beta_1\cdots\beta_t\) is a \emph{complete} factorization,
	every \(i \in X = \{1, \dots, n\}\) occurs in exactly one \(\beta_\lambda\); hence
	\(\alpha(i)\) is defined for every \(i \in X\) and \(\alpha\colon X \to X\) is a
	single-valued function. Since every \(j \in X\) occurs in some \(\sigma_\lambda\),
	\(\alpha\) is surjective, hence a bijection, and so \(\alpha \in S_n\). By
	\autoref{prop:conj-same-structure}, the \(\lambda\)th cycle of
	\(\alpha\gamma\alpha^{-1}\) is
	\(\bigl(\alpha(i_1^\lambda)\ \dots\ \alpha(i_{r(\lambda)}^\lambda)\bigr) = \sigma_\lambda\).
	Therefore \(\alpha\gamma\alpha^{-1} = \gamma'\).
\end{proof}

\begin{eg}
	If \(\gamma = (1\ 2\ 3)(4\ 5)(6)\) and \(\gamma' = (2\ 5\ 6)(3\ 1)(4)\), then
	\(\gamma' = \alpha\gamma\alpha^{-1}\) where
	\[
		\alpha = \begin{pmatrix} 1&2&3&4&5&6 \\ 2&5&6&3&1&4 \end{pmatrix}
		= (1\ 2\ 5)(3\ 6\ 4).
	\]
	Note that there are other choices for \(\alpha\) as well.
\end{eg}

\subsection{Parity}

\begin{proposition}\label{prop:product-of-transpositions}
	If \(n \geq 2\), then every \(\alpha \in S_n\) is a product of transpositions.
\end{proposition}

\begin{proof}
	Of course \((1) = (1\ 2)(1\ 2)\) is such a product, as is every transposition:
	\((i\ j) = (i\ j)(1\ 2)(1\ 2)\). By \autoref{prop:disjoint-cycles} it suffices to
	factor an \(r\)-cycle \(\beta\). If \(r = 1\), then \(\beta = (1\ 2)(1\ 2)\). If
	\(r \geq 2\), then
	\[
		\beta = (1\ 2\ \dots\ r) = (1\ r)(1\ r-1)\cdots(1\ 3)(1\ 2).
	\]
	One checks this by evaluating each side: the left side sends \(1 \mapsto 2\), and each of
	\((1\ r), (1\ r-1), \dots, (1\ 3)\) fixes \(2\), so the right side also sends
	\(1 \mapsto 2\); and so on.
\end{proof}

\begin{remark}[Transposition factorizations are far from unique]
	Such a factorization is not as nice as the factorization into disjoint cycles. The
	transpositions occurring need not commute:
	\((1\ 2\ 3) = (1\ 3)(1\ 2) \neq (1\ 2)(1\ 3)\); and neither the factors nor their
	number is determined. For example, in \(S_4\),
	\begin{align*}
		(1\ 2\ 3) & = (1\ 3)(1\ 2)                              \\
		          & = (2\ 3)(1\ 3)                              \\
		          & = (1\ 3)(4\ 2)(1\ 2)(1\ 4)                  \\
		          & = (1\ 3)(4\ 2)(1\ 2)(1\ 4)(2\ 3)(2\ 3).
	\end{align*}
	The only invariant that survives is the \emph{parity} of the number of factors, and
	proving this is the goal of the rest of the section.
\end{remark}

\begin{eg}[The \(15\)-puzzle]\label{eg:15-puzzle}
	The \vocab{\(15\)-puzzle} consists of a starting position, which is a \(4 \times 4\)
	array of the numbers \(1\) through \(15\) together with a symbol \(\#\) (a blank), and
	\vocab{simple moves}. Consider the starting position
	\[
		\begin{array}{|c|c|c|c|}
			\hline
			3  & 15 & 4  & 12 \\ \hline
			10 & 11 & 1  & 8  \\ \hline
			2  & 5  & 13 & 9  \\ \hline
			6  & 7  & 14 & \# \\ \hline
		\end{array}
	\]
	A simple move interchanges the blank with an adjacent symbol; here there are two
	beginning moves, interchanging \(\#\) with \(14\) or with \(9\). One \emph{wins} if,
	after a sequence of simple moves, the starting position is transformed into the standard
	array \(1, 2, \dots, 15, \#\).

	To analyze the game, note that the given array is really a permutation
	\(\alpha \in S_{16}\) of \(\{1, 2, \dots, 15, \#\}\): if the \emph{squares} are labelled
	\(1\) through \(15, \#\), then \(\alpha(i)\) is the symbol occupying the \(i\)th square.
	The position above is
	\[
		\alpha =
		\setlength{\arraycolsep}{2.5pt}
		\left(\begin{array}{*{16}{c}}
				1 & 2  & 3 & 4  & 5  & 6  & 7 & 8 & 9 & 10 & 11 & 12 & 13 & 14 & 15 & \# \\
				3 & 15 & 4 & 12 & 10 & 11 & 1 & 8 & 2 & 5  & 13 & 9  & 6  & 7  & 14 & \#
			\end{array}\right).
	\]
	Each simple move is a transposition \(\tau\) that moves \(\#\), and performing it from
	the position \(\beta\) yields the position \(\tau\beta\). Hence to win one needs special
	transpositions \(\tau_1, \dots, \tau_m\) with
	\[
		\tau_m \cdots \tau_2\tau_1\alpha = (1).
	\]
	Some starting positions can be won and others cannot, as we shall see in
	\autoref{eg:15-puzzle-solved}.
\end{eg}

\begin{lemma}\label{lem:split-merge}
	If \(k, \ell \geq 0\) and the letters \(a, b, c_i, d_j\) are all distinct, then
	\[
		(a\ b)(a\ c_1 \dots c_k\ b\ d_1 \dots d_\ell) = (a\ c_1 \dots c_k)(b\ d_1 \dots d_\ell)
	\]
	and
	\[
		(a\ b)(a\ c_1 \dots c_k)(b\ d_1 \dots d_\ell) = (a\ c_1 \dots c_k\ b\ d_1 \dots d_\ell).
	\]
\end{lemma}

\begin{proof}
	The left side of the first equation sends
	\[
		a \mapsto c_1 \mapsto c_1; \quad
		c_i \mapsto c_{i+1} \mapsto c_{i+1} \ (i < k); \quad
		c_k \mapsto b \mapsto a;
	\]
	\[
		b \mapsto d_1 \mapsto d_1; \quad
		d_j \mapsto d_{j+1} \mapsto d_{j+1} \ (j < \ell); \quad
		d_\ell \mapsto a \mapsto b.
	\]
	Evaluating the right side gives the same values on \(a\), \(b\), and all \(c_i, d_j\),
	and each side fixes every other number, so the two sides are equal.

	For the second equation, reverse the first and multiply both sides on the left by
	\((a\ b)\), using \((a\ b)(a\ b) = (1)\).
\end{proof}

An illustration: \((1\ 2)(1\ 3\ 4\ 2\ 5\ 6\ 7) = (1\ 3\ 4)(2\ 5\ 6\ 7)\).

\begin{definition}[Signum]
	If \(\alpha \in S_n\) and \(\alpha = \beta_1 \cdots \beta_t\) is a \emph{complete}
	factorization into disjoint cycles, then the \vocab{signum} of \(\alpha\) is
	\[
		\sgn(\alpha) = (-1)^{n-t}.
	\]
\end{definition}

\autoref{thm:unique-factorization} shows that \(\sgn\) is a well-defined function, for
\(t\) is determined by \(\alpha\). If \(\varepsilon\) is a \(1\)-cycle, then \(t = n\) and
\(\sgn(\varepsilon) = 1\). If \(\tau\) is a transposition, then it moves two numbers and
fixes the other \(n - 2\), so \(t = 1 + (n-2) = n-1\) and
\(\sgn(\tau) = (-1)^{n-(n-1)} = -1\).

\begin{lemma}\label{lem:sgn-transposition}
	If \(\alpha, \tau \in S_n\) with \(\tau\) a transposition, then
	\(\sgn(\tau\alpha) = -\sgn(\alpha)\).
\end{lemma}

\begin{proof}
	Let \(\alpha = \beta_1 \cdots \beta_t\) be a complete factorization and let
	\(\tau = (a\ b)\).

	If \(a\) and \(b\) occur in the same \(\beta\), say \(\beta_1\), then
	\(\beta_1 = (a\ c_1 \dots c_k\ b\ d_1 \dots d_\ell)\) with \(k, \ell \geq 0\), and
	\autoref{lem:split-merge} gives
	\(\tau\beta_1 = (a\ c_1 \dots c_k)(b\ d_1 \dots d_\ell)\). This exhibits a complete
	factorization of \(\tau\alpha = (\tau\beta_1)\beta_2\cdots\beta_t\) with \(t+1\) cycles,
	for \(\beta_1\) has split in two. Therefore
	\(\sgn(\tau\alpha) = (-1)^{n-(t+1)} = -\sgn(\alpha)\).

	Otherwise \(a\) and \(b\) occur in different cycles, say
	\(\beta_1 = (a\ c_1 \dots c_k)\) and \(\beta_2 = (b\ d_1 \dots d_\ell)\). Then
	\(\tau\alpha = (\tau\beta_1\beta_2)\beta_3\cdots\beta_t\) and
	\autoref{lem:split-merge} gives
	\(\tau\beta_1\beta_2 = (a\ c_1 \dots c_k\ b\ d_1 \dots d_\ell)\). So \(\tau\alpha\) has a
	complete factorization with \(t-1\) cycles, and
	\(\sgn(\tau\alpha) = (-1)^{n-(t-1)} = -\sgn(\alpha)\).
\end{proof}

\begin{theorem}[\(\sgn\) is multiplicative]\label{thm:sgn-multiplicative}
	For all \(\alpha, \beta \in S_n\),
	\[
		\sgn(\alpha\beta) = \sgn(\alpha)\sgn(\beta).
	\]
\end{theorem}

\begin{proof}
	By \autoref{prop:product-of-transpositions} we may write
	\(\alpha = \tau_1 \cdots \tau_m\) as a product of transpositions; induct on \(m\). The
	base step \(m = 1\) is \autoref{lem:sgn-transposition}. If \(m > 1\), the inductive
	hypothesis applies to \(\tau_2\cdots\tau_m\), and
	\begin{align*}
		\sgn(\alpha\beta) & = \sgn(\tau_1\cdots\tau_m\beta)                    \\
		                  & = -\sgn(\tau_2\cdots\tau_m\beta)                   && \text{(\autoref{lem:sgn-transposition})} \\
		                  & = -\sgn(\tau_2\cdots\tau_m)\sgn(\beta)             && \text{(induction)}                         \\
		                  & = \sgn(\tau_1\cdots\tau_m)\sgn(\beta)              && \text{(\autoref{lem:sgn-transposition})} \\
		                  & = \sgn(\alpha)\sgn(\beta). &&\qedhere
	\end{align*}
\end{proof}

It follows by induction on \(k \geq 2\) that
\(\sgn(\alpha_1\alpha_2\cdots\alpha_k) = \sgn(\alpha_1)\sgn(\alpha_2)\cdots\sgn(\alpha_k)\).

\begin{definition}[Even, odd, parity]
	A permutation \(\alpha \in S_n\) is \vocab{even} if \(\sgn(\alpha) = 1\), and
	\vocab{odd} if \(\sgn(\alpha) = -1\). We say \(\alpha\) and \(\beta\) have the same
	\vocab{parity} if both are even or both are odd.
\end{definition}

\begin{theorem}\label{thm:parity}
	\begin{enumerate}[(i)]
		\item Let \(\alpha \in S_n\). If \(\alpha\) is even, then \(\alpha\) is a product of
		      an even number of transpositions; if \(\alpha\) is odd, then \(\alpha\) is a
		      product of an odd number of transpositions.
		\item If \(\alpha = \tau_1\cdots\tau_q = \tau'_1\cdots\tau'_p\) are two
		      factorizations into transpositions, then \(q\) and \(p\) have the same parity.
	\end{enumerate}
\end{theorem}

\begin{proof}
	(i) If \(\alpha = \tau_1\cdots\tau_q\), then \autoref{thm:sgn-multiplicative} gives
	\(\sgn(\alpha) = \sgn(\tau_1)\cdots\sgn(\tau_q) = (-1)^q\), since every transposition is
	odd. Hence \(\sgn(\alpha) = 1\) forces \(q\) even, and \(\sgn(\alpha) = -1\) forces
	\(q\) odd.

	(ii) If there were two factorizations, one into an odd and one into an even number of
	transpositions, then \(\sgn(\alpha)\) would have two different values.
\end{proof}

\begin{corollary}\label{cor:parity-product}
	Let \(\alpha, \beta \in S_n\). If \(\alpha\) and \(\beta\) have the same parity, then
	\(\alpha\beta\) is even; if they have distinct parity, then \(\alpha\beta\) is odd.
\end{corollary}

\begin{proof}
	If \(\sgn(\alpha) = (-1)^q\) and \(\sgn(\beta) = (-1)^p\), then
	\autoref{thm:sgn-multiplicative} gives \(\sgn(\alpha\beta) = (-1)^{q+p}\).
\end{proof}

\begin{eg}[The \(15\)-puzzle, resolved]\label{eg:15-puzzle-solved}
	One can show that if \(\alpha \in S_{16}\) is the starting position of the \(15\)-puzzle,
	then the game can be won if and only if \(\alpha\) is an \emph{even} permutation fixing
	\(\#\). One direction is clear. The blank \(\#\) starts in position \(16\); each simple
	move takes \(\#\) up, down, left, or right, so the total number of moves is
	\(m = u + d + l + r\). If \(\#\) is to return home, then \(u = d\) and \(l = r\), so
	\(m = 2u + 2r\) is even. That is, if \(\tau_m\cdots\tau_1\alpha = (1)\), then \(m\) is
	even; hence \(\alpha = \tau_1\cdots\tau_m\) (because \(\tau^{-1} = \tau\) for every
	transposition), and so \(\alpha\) is even.

	For the starting position of \autoref{eg:15-puzzle}, the cycle notation is
	\[
		\alpha = (1\ 3\ 4\ 12\ 9\ 2\ 15\ 14\ 7)(5\ 10)(6\ 11\ 13)(8)(\#),
	\]
	so \(t = 5\) and \(\sgn(\alpha) = (-1)^{16-5} = -1\). Therefore \(\alpha\) is odd, and
	the game starting with \(\alpha\) \emph{cannot} be won.
\end{eg}

\section{Groups}

Generalizations of the quadratic formula for cubics and quartics were discovered in the
early 1500s. Over the next three centuries many tried to find analogous formulas for
higher degree, but in 1824 N.~H.~Abel (1802--1829) proved that no such formula exists for
the general polynomial of degree \(5\). In 1831 E.~Galois (1811--1832) solved the problem
completely, determining precisely which polynomials admit such a formula; his fundamental
idea was the invention of the notion of \emph{group}. Since Galois's time groups have
arisen throughout mathematics, for they are also the way to describe \emph{symmetry}.

\begin{intuition}
	The essence of a ``product'' is that two things are combined to form a third thing of
	the same kind: multiplication, addition and subtraction combine two numbers to give a
	number; composition combines two permutations to give a permutation.
\end{intuition}

\begin{definition}[Binary operation]
	A \vocab{(binary) operation} on a set \(G\) is a function
	\(* \colon G \times G \to G\).
\end{definition}

In more detail, an operation assigns an element \(*(x,y) \in G\) to each \emph{ordered}
pair \((x,y)\) of elements of \(G\); it is more natural to write \(x * y\). Thus
composition of functions is the operation \((f,g) \mapsto g \circ f\), while
multiplication, addition and subtraction are \((x,y) \mapsto xy\),
\((x,y) \mapsto x+y\) and \((x,y) \mapsto x-y\). Composition and subtraction show why we
want \emph{ordered} pairs, for \(x * y\) and \(y * x\) may differ.

\begin{note}[Law of substitution]
	An operation is, as any function, single-valued. Stated explicitly, this is called the
	\vocab{law of substitution}: if \(x = x'\) and \(y = y'\), then \(x*y = x'*y'\).
\end{note}

\begin{definition}[Group]\label{def:group}
	A \vocab{group} is a set \(G\) equipped with an operation \(*\) and a special element
	\(e \in G\), called the \vocab{identity}, such that
	\begin{enumerate}[(i)]
		\item the \vocab{associative law} holds: \(a*(b*c) = (a*b)*c\) for all
		      \(a,b,c \in G\);
		\item \(e * a = a\) for all \(a \in G\);
		\item for every \(a \in G\) there is \(a' \in G\) with \(a' * a = e\).
	\end{enumerate}
\end{definition}

\begin{remark}
	Only a \emph{left} identity and \emph{left} inverses are assumed. The two-sided
	statements \(a * e = a\) and \(a * a' = e\) are theorems, proved in
	\autoref{prop:group-basics}.
\end{remark}

The set \(S_X\) of all permutations of a set \(X\), with composition as operation and
\(1_X\) as identity, is a group: the symmetric group on \(X\).

\begin{intuition}[Why abstraction pays]
	We are now at the precise point where algebra becomes \emph{abstract} algebra. In
	contrast to the concrete group \(S_n\), we will prove general results about groups
	without specifying either their elements or their operation. Theorems then apply to
	\emph{many} different groups at once, and it is often simpler to work abstractly even
	when dealing with one particular concrete group (see \autoref{eg:finite-order}).
\end{intuition}

\begin{definition}[Abelian group]
	A group \(G\) is \vocab{abelian} if it satisfies the \vocab{commutative law}:
	\(x*y = y*x\) for all \(x, y \in G\).
\end{definition}

The name honours Abel, who proved in 1827 a theorem equivalent to the existence of a
formula for the roots of a polynomial whose Galois group is commutative. The groups
\(S_n\) for \(n \geq 3\) are \emph{not} abelian, because
\((1\ 2)(1\ 3) = (1\ 3\ 2)\) while \((1\ 3)(1\ 2) = (1\ 2\ 3)\).

\subsection{First Consequences of the Axioms}

Associativity says \(a*(b*c) = (a*b)*c\), so it is unambiguous to write \(a*b*c\). Not
all operations are associative: if \(c \neq 0\), then
\(a - (b-c) \neq (a-b) - c\). The next lemma shows that some associativity carries over to
products of four factors.

\begin{lemma}\label{lem:assoc4}
	If \(*\) is an associative operation on a set \(G\), then
	\[
		(a*b)*(c*d) = \bigl[a*(b*c)\bigr]*d
	\]
	for all \(a,b,c,d \in G\).
\end{lemma}

\begin{proof}
	Write \(g = a*b\). Then
	\((a*b)*(c*d) = g*(c*d) = (g*c)*d = \bigl[(a*b)*c\bigr]*d = \bigl[a*(b*c)\bigr]*d\).
\end{proof}

\begin{lemma}\label{lem:idempotent}
	If \(G\) is a group and \(a \in G\) satisfies \(a*a = a\), then \(a = e\).
\end{lemma}

\begin{proof}
	There is \(a' \in G\) with \(a'*a = e\). Multiplying both sides of \(a*a = a\) on the
	left by \(a'\) gives \(a'*(a*a) = a'*a = e\); but
	\(a'*(a*a) = (a'*a)*a = e*a = a\). Hence \(a = e\).
\end{proof}

\begin{proposition}\label{prop:group-basics}
	Let \(G\) be a group with operation \(*\) and identity \(e\).
	\begin{enumerate}[(i)]
		\item \(a*a' = e\) for all \(a \in G\).
		\item \(a*e = a\) for all \(a \in G\).
		\item If \(e' \in G\) satisfies \(e'*a = a\) for all \(a \in G\), then \(e' = e\);
		      that is, the identity is unique.
		\item Let \(a \in G\). If \(b \in G\) satisfies \(b*a = e\), then \(b = a'\); that
		      is, inverses are unique.
	\end{enumerate}
\end{proposition}

\begin{proof}
	(i) By \autoref{lem:assoc4},
	\[
		(a*a')*(a*a') = \bigl[a*(a'*a)\bigr]*a' = (a*e)*a' = a*(e*a') = a*a',
	\]
	and \autoref{lem:idempotent} gives \(a*a' = e\).

	(ii) Using part (i),
	\(a*e = a*(a'*a) = (a*a')*a = e*a = a\).

	(iii) If \(e'*a = a\) for all \(a \in G\), then in particular \(e'*e' = e'\), so
	\(e' = e\) by \autoref{lem:idempotent}.

	(iv) By part (i), \(a*a' = e\). Hence
	\(b = b*e = b*(a*a') = (b*a)*a' = e*a' = a'\).
\end{proof}

\begin{definition}[Inverse]
	If \(G\) is a group and \(a \in G\), then the unique element \(a' \in G\) with
	\(a'*a = e\) is called the \vocab{inverse} of \(a\), denoted by \(a^{-1}\).
\end{definition}

\begin{lemma}\label{lem:group-cancellation}
	Let \(G\) be a group.
	\begin{enumerate}[(i)]
		\item The \vocab{cancellation laws} hold: if \(a,b,x \in G\) and either
		      \(x*a = x*b\) or \(a*x = b*x\), then \(a = b\).
		\item \((a^{-1})^{-1} = a\) for all \(a \in G\).
		\item If \(a, b \in G\), then \((a*b)^{-1} = b^{-1}*a^{-1}\). More generally, for
		      all \(n \geq 2\),
		      \[
			      (a_1 * a_2 * \cdots * a_n)^{-1} = a_n^{-1} * \cdots * a_2^{-1} * a_1^{-1}.
		      \]
	\end{enumerate}
\end{lemma}

\begin{proof}
	(i)
	\[
		a = e*a = (x^{-1}*x)*a = x^{-1}*(x*a) = x^{-1}*(x*b) = (x^{-1}*x)*b = e*b = b.
	\]
	A similar proof, using \(x*x^{-1} = e\), works when \(x\) is on the right.

	(ii) By \autoref{prop:group-basics}(i), \(a*a^{-1} = e\). Uniqueness of
	inverses, \autoref{prop:group-basics}(iv), says that \((a^{-1})^{-1}\) is the
	unique \(x \in G\) with \(x*a^{-1} = e\); therefore \((a^{-1})^{-1} = a\).

	(iii) By \autoref{lem:assoc4},
	\[
		(a*b)*(b^{-1}*a^{-1}) = \bigl[a*(b*b^{-1})\bigr]*a^{-1} = (a*e)*a^{-1} = a*a^{-1} = e,
	\]
	so \((a*b)^{-1} = b^{-1}*a^{-1}\) by \autoref{prop:group-basics}(iv). The
	general statement follows by induction on \(n\).
\end{proof}

\begin{notation}
	From now on we usually write \(ab\) for \(a*b\) and \(1\) for the identity. When a
	group is abelian we often use \emph{additive} notation instead: an \vocab{additive
		group} is a set \(G\) with an operation \(+\) and an element \(0 \in G\) such that
	\(a + (b+c) = (a+b)+c\), \(0 + a = a\), and for every \(a\) there is \(-a \in G\) with
	\((-a) + a = 0\). The inverse of \(a\) is written \(-a\) rather than \(a^{-1}\).
\end{notation}

\subsection{Examples of Groups}

\begin{eg}[A first list of groups]\label{eg:group-list}
	\begin{enumerate}[(i)]
		\item \(S_X\), the set of all permutations of a set \(X\), is a group under
		      composition; in particular \(S_n\) is a group.
		\item \(\mathbb{Z}\) is an additive abelian group with \(a*b = a+b\), identity
		      \(e = 0\), and inverse \(-n\) of \(n\). Similarly \(\mathbb{Q}\),
		      \(\mathbb{R}\) and \(\mathbb{C}\) are additive abelian groups.
		\item \(\mathbb{Q}^\times\), the set of all nonzero rationals, is an abelian group
		      under ordinary multiplication, with identity \(1\) and inverse \(1/r\) of
		      \(r\). Similarly \(\mathbb{R}^\times\) is a multiplicative abelian group.
		      Note that \(\mathbb{Z}^\times\) is \emph{not} a group, for no element other
		      than \(\pm 1\) has a multiplicative inverse in \(\mathbb{Z}^\times\).
		\item The nonzero complex numbers \(\mathbb{C}^\times\) form an abelian group under
		      multiplication. For inverses, if \(z = a + ib\) with \(a, b \in \mathbb{R}\),
		      let \(\overline{z} = a - ib\) be its complex conjugate; then
		      \(z\overline{z} = a^2 + b^2\), so \(z \neq 0\) if and only if
		      \(z\overline{z} \neq 0\), and for \(z \neq 0\),
		      \[
			      z^{-1} = \overline{z}/z\overline{z} = (a/z\overline{z}) - (b/z\overline{z})i.
		      \]
		\item The \vocab{circle group} is
		      \(S^1 = \{z \in \mathbb{C} : \abs{z} = 1\}\) under multiplication of complex
		      numbers; the identity is \(1\), and the inverse of a number of modulus \(1\)
		      is its complex conjugate. Although \(S^1\) is abelian we still write it
		      multiplicatively.
		\item For a positive integer \(n\), let
		      \(\Gamma_n = \{\zeta^k : 0 \leq k < n\}\) be the set of all \(n\)th roots of
		      unity, where \(\zeta = e^{2\pi i/n} = \cos(2\pi/n) + i\sin(2\pi/n)\). De
		      Moivre's theorem shows \(\Gamma_n\) is an abelian group under multiplication,
		      the inverse of a root of unity being its complex conjugate.
		\item The plane \(\mathbb{R} \times \mathbb{R}\) is an additive abelian group under
		      vector addition: \((x,y) + (x',y') = (x+x', y+y')\). The identity is the
		      origin \(O = (0,0)\) and the inverse of \(\vb{v} = (x,y)\) is
		      \(-\vb{v} = (-x,-y)\).
		\item The \vocab{parity group} \(\mathbb{P}\) has two elements, the words ``even''
		      and ``odd'', with
		      \[
			      \text{even} + \text{even} = \text{even} = \text{odd} + \text{odd},
			      \qquad
			      \text{even} + \text{odd} = \text{odd} = \text{odd} + \text{even}.
		      \]
		\item The \vocab{Boolean group} \(\mathcal{B}(X)\) of a set \(X\) (named after the
		      logician G.~Boole, 1815--1864) is the family of all subsets of \(X\) with
		      addition given by symmetric difference
		      \(A + B = (A - B) \cup (B - A)\). Symmetric difference is commutative, the
		      identity is \(\varnothing\), and every element is its own inverse, since
		      \(A + A = \varnothing\). Thus \(\mathcal{B}(X)\) is abelian.
	\end{enumerate}
\end{eg}

\begin{eg}[Matrix groups and affine maps]\label{eg:matrix-groups}
	\begin{enumerate}[(i)]
		\item A \((2\times 2\) real\()\) matrix is
		      \(A = \begin{psmallmatrix} a & c \\ b & d \end{psmallmatrix}\) with
		      \(a,b,c,d \in \mathbb{R}\), and the product is
		      \[
			      AB =
			      \begin{pmatrix} a & c \\ b & d \end{pmatrix}
			      \begin{pmatrix} w & y \\ x & z \end{pmatrix}
			      =
			      \begin{pmatrix} aw + cx & ay + cz \\ bw + dx & by + dz \end{pmatrix};
		      \]
		      each entry of \(AB\) is a dot product of a row of \(A\) with a column of
		      \(B\). The \vocab{determinant} of \(A\) is \(\det(A) = ad - bc\), and \(A\) is
		      \vocab{nonsingular} if \(\det(A) \neq 0\). Since
		      \(\det(AB) = \det(A)\det(B)\), a product of nonsingular matrices is
		      nonsingular. The set \(GL(2,\mathbb{R})\) of all nonsingular matrices, under
		      matrix multiplication, is a \emph{nonabelian} group, the \(2\times 2\) real
		      \vocab{general linear group}: the identity is
		      \(I = \begin{psmallmatrix} 1&0 \\ 0&1\end{psmallmatrix}\) and
		      \[
			      A^{-1} = \begin{pmatrix} d/\Delta & -c/\Delta \\ -b/\Delta & a/\Delta \end{pmatrix},
			      \qquad \Delta = ad - bc = \det(A).
		      \]
		\item Allowing entries in \(\mathbb{Q}\) or \(\mathbb{C}\) gives
		      \(GL(2,\mathbb{Q})\) and \(GL(2,\mathbb{C})\); one may even allow entries in
		      \(\mathbb{Z}\), in which case \(GL(2,\mathbb{Z})\) is defined to be the set of
		      such matrices with determinant \(\pm 1\) (so that \(A^{-1}\) also has integer
		      entries). More generally, all nonsingular \(n \times n\) matrices form a group
		      \(GL(n, k)\).
		\item All \vocab{special orthogonal} matrices, that is, all matrices
		      \(A = \begin{psmallmatrix} \cos\alpha & -\sin\alpha \\
				      \sin\alpha & \cos\alpha \end{psmallmatrix}\), form the \(2 \times 2\)
		      \vocab{special orthogonal group} \(SO(2,\mathbb{R})\). Indeed
		      \[
			      \begin{pmatrix} \cos\alpha & -\sin\alpha \\ \sin\alpha & \cos\alpha \end{pmatrix}
			      \begin{pmatrix} \cos\beta  & -\sin\beta  \\ \sin\beta  & \cos\beta  \end{pmatrix}
			      =
			      \begin{pmatrix} \cos(\alpha+\beta) & -\sin(\alpha+\beta) \\
				      \sin(\alpha+\beta) & \cos(\alpha+\beta) \end{pmatrix}
		      \]
		      by the addition theorems for sine and cosine; this calculation also shows that
		      \(SO(2,\mathbb{R})\) is abelian. The identity matrix is special orthogonal, and
		      so is the inverse of a special orthogonal matrix. We shall see that
		      \(SO(2,\mathbb{R})\) is a disguised version of the circle group \(S^1\), and
		      that it consists of all rotations of the plane about the origin.
		\item The \vocab{affine group} \(\mathrm{Aff}(1,\mathbb{R})\) consists of all
		      \vocab{affine maps} \(f_{a,b}\colon \mathbb{R}\to\mathbb{R}\),
		      \(f_{a,b}(x) = ax + b\), with \(a, b\) fixed reals and \(a \neq 0\). It is a
		      group under composition: if \(f_{c,d}(x) = cx+d\), then
		      \[
			      f_{a,b}f_{c,d}(x) = a(cx+d) + b = acx + (ad+b) = f_{ac,\,ad+b}(x),
		      \]
		      and \(ac \neq 0\); the identity is \(1_{\mathbb{R}} = f_{1,0}\), and the
		      inverse of \(f_{a,b}\) is \(f_{a^{-1},\,-a^{-1}b}\). The composition is
		      reminiscent of matrix multiplication:
		      \[
			      \begin{pmatrix} a & b \\ 0 & 1 \end{pmatrix}
			      \begin{pmatrix} c & d \\ 0 & 1 \end{pmatrix}
			      =
			      \begin{pmatrix} ac & ad+b \\ 0 & 1 \end{pmatrix}.
		      \]
	\end{enumerate}
\end{eg}

\subsection{Generalized Associativity and Powers}

\begin{definition}[\(n\)-expression, ultimate product]
	An \vocab{\(n\)-expression} is an \(n\)-tuple \((a_1, \dots, a_n)\) in the product of
	\(n\) copies of \(G\). It yields elements of \(G\) as follows: choose two adjacent
	entries of the tuple, multiply them, and
	obtain an \((n-1)\)-expression; repeat until a \(2\)-expression \((W,X)\) is reached,
	and then form \(WX \in G\). Such an element \(WX\) is called an \vocab{ultimate
		product} derived from the expression. The expression \vocab{needs no parentheses} if
	all ultimate products derived from it are equal.
\end{definition}

For example, from \((a,b,c,d)\) one can reach \([(ab)c]d\), \((ab)(cd)\), \(a[(bc)d]\),
and so on. To say that an operation is associative is exactly to say that the two ultimate
products from every \(3\)-expression agree; it is not obvious that all ultimate products
from a longer expression agree.

\begin{theorem}[Generalized associativity]\label{thm:gen-assoc}
	If \(n \geq 3\), then every \(n\)-expression \((a_1, \dots, a_n)\) in a group \(G\)
	needs no parentheses.
\end{theorem}

\begin{remark}
	Neither the identity element nor inverses are used in the proof; the hypothesis can be
	weakened to \(G\) being a \vocab{semigroup}, that is, a nonempty set with an associative
	binary operation.
\end{remark}

\begin{proof}
	Induct on \(n\) (second form). The base step \(n = 3\) is associativity. For the
	inductive step, consider two \(2\)-expressions obtained from \((a_1, \dots, a_n)\):
	\[
		(W, X) = (a_1\cdots a_i,\ a_{i+1}\cdots a_n)
		\quad\text{and}\quad
		(Y, Z) = (a_1\cdots a_j,\ a_{j+1}\cdots a_n).
	\]
	We must prove \(WX = YZ\). By induction each of \(W, X, Y, Z\) is the one and only
	ultimate product of an \(m\)-expression with \(m < n\). Without loss of generality
	\(i \leq j\). If \(i = j\), then \(W = Y\) and \(X = Z\), so \(WX = YZ\).

	Assume \(i < j\). Let \(A\) be the ultimate product of \((a_1, \dots, a_i)\), let \(B\)
	be that of \((a_{i+1}, \dots, a_j)\), and let \(C\) be that of
	\((a_{j+1}, \dots, a_n)\); these are unambiguous by the inductive hypothesis. Then
	\(W = A\), \(Z = C\), \(X = BC\) and \(Y = AB\), so \(WX = A(BC)\) and
	\(YZ = (AB)C\). The base step \(n = 3\) gives \(WX = YZ\).
\end{proof}

\begin{definition}[Powers]
	If \(G\) is a group and \(a \in G\), define \(a^n\) for \(n \geq 1\) inductively by
	\(a^1 = a\) and \(a^{n+1} = a a^n\); define \(a^0 = 1\) and, for \(n > 0\),
	\(a^{-n} = (a^{-1})^n\).
\end{definition}

One shows \((a^{-1})^n = (a^n)^{-1}\); this is a special case of
\autoref{lem:group-cancellation}(iii). There is a hidden complication: we defined
\(a^3 = a a^2 = a(aa)\), but \((aa)a = a^2 a\) is another reasonable contender.
Associativity makes them equal, and generalized associativity shows that all powers are
unambiguously defined.

\begin{corollary}\label{cor:powers}
	If \(G\) is a group, \(a \in G\), and \(m, n \geq 1\), then
	\[
		a^{m+n} = a^m a^n \quad\text{and}\quad (a^m)^n = a^{mn}.
	\]
\end{corollary}

\begin{proof}
	Both \(a^{m+n}\) and \(a^m a^n\) are ultimate products of the expression having \(m+n\)
	factors each equal to \(a\); both \((a^m)^n\) and \(a^{mn}\) arise from the expression
	having \(mn\) factors each equal to \(a\). Apply \autoref{thm:gen-assoc}.
\end{proof}

In particular any two powers of the same element commute:
\(a^m a^n = a^{m+n} = a^{n+m} = a^n a^m\).

\begin{proposition}[Laws of exponents]\label{prop:exponents}
	Let \(G\) be a group, \(a, b \in G\), and let \(m, n\) be (not necessarily positive)
	integers.
	\begin{enumerate}[(i)]
		\item If \(a\) and \(b\) commute, then \((ab)^n = a^n b^n\).
		\item \((a^n)^m = a^{mn}\).
		\item \(a^m a^n = a^{m+n}\).
	\end{enumerate}
\end{proposition}

\begin{proof}
	Exercises for the reader.
\end{proof}

\begin{notation}
	If the operation is \(+\), it is more natural to write \(na\) for
	\(a + a + \cdots + a\). In additive notation \autoref{prop:exponents} reads
	\[
		n(a+b) = na + nb, \qquad m(na) = (mn)a, \qquad ma + na = (m+n)a.
	\]
\end{notation}

\subsection{The Order of an Element}

\begin{eg}[Shuffling a deck]\label{eg:finite-order}
	Suppose a deck of cards is shuffled, changing the order from
	\(1, 2, 3, 4, \dots, 52\) to \(2, 1, 4, 3, \dots, 52, 51\). Shuffling again the same way
	restores the original order. In fact for \emph{any} permutation \(\alpha\) of the
	\(52\) cards, repeating \(\alpha\) sufficiently often restores the deck. One proof uses
	our knowledge of permutations: write \(\alpha = \beta_1\cdots\beta_t\) with \(\beta_i\)
	a disjoint \(r_i\)-cycle. Then \(\beta_i^{r_i} = (1)\) by
	\autoref{lem:cycle-order}, so \(\beta_i^k = (1)\) for \(k = r_1\cdots r_t\); since
	disjoint cycles commute, \autoref{prop:exponents}(i) gives
	\[
		\alpha^k = (\beta_1\cdots\beta_t)^k = \beta_1^k \cdots \beta_t^k = (1).
	\]

	Here is a more general result with a \emph{simpler} proof (abstract algebra can be
	easier than algebra): if \(G\) is a finite group and \(a \in G\), then \(a^k = 1\) for
	some \(k \geq 1\). Indeed, consider \(1, a, a^2, \dots\). Since \(G\) is finite there
	must be a repetition, say \(a^m = a^n\) with \(m > n\); then
	\(1 = a^m a^{-n} = a^{m-n}\).
\end{eg}

\begin{definition}[Order of an element]
	Let \(G\) be a group and \(a \in G\). If \(a^k = 1\) for some \(k \geq 1\), then the
	smallest such exponent \(k \geq 1\) is called the \vocab{order} of \(a\); if no such
	power exists, \(a\) is said to have \vocab{infinite order}.
\end{definition}

Every element of a finite group has finite order, by \autoref{eg:finite-order}. In any
group the identity is the only element of order \(1\); an element has order \(2\) if and
only if it is not the identity and equals its own inverse. The matrix
\(A = \begin{psmallmatrix} 1&1 \\ 0&1\end{psmallmatrix} \in GL(2,\mathbb{R})\) has infinite
order, since \(A^k = \begin{psmallmatrix} 1&k \\ 0&1\end{psmallmatrix} \neq I\) for all
\(k \geq 1\).

\begin{lemma}\label{lem:order-divides}
	Let \(G\) be a group and let \(a \in G\) have finite order \(k\). If \(a^n = 1\), then
	\(k \mid n\). In fact \(\{n \in \mathbb{Z} : a^n = 1\}\) is the set of all multiples of
	\(k\).
\end{lemma}

\begin{proof}
	Let \(I = \{n \in \mathbb{Z} : a^n = 1\}\). Then \(0 \in I\) because \(a^0 = 1\); if
	\(n, m \in I\), then \(a^{n-m} = a^n a^{-m} = 1\), so \(n - m \in I\); and if
	\(n \in I\), \(q \in \mathbb{Z}\), then \(a^{qn} = (a^n)^q = 1\), so \(qn \in I\).
	Therefore \(I\) consists of all the multiples of its smallest positive element, which by
	definition is the order \(k\) of \(a\). Hence \(a^n = 1\) implies \(k \mid n\).
\end{proof}

\begin{proposition}[Order of a permutation]\label{prop:order-permutation}
	Let \(\alpha \in S_n\).
	\begin{enumerate}[(i)]
		\item If \(\alpha\) is an \(r\)-cycle, then \(\alpha\) has order \(r\).
		\item If \(\alpha = \beta_1\cdots\beta_t\) is a product of disjoint \(r_i\)-cycles
		      \(\beta_i\), then \(\alpha\) has order
		      \(m = \operatorname{lcm}(r_1, \dots, r_t)\).
		\item If \(p\) is a prime, then \(\alpha\) has order \(p\) if and only if it is a
		      \(p\)-cycle or a product of disjoint \(p\)-cycles.
	\end{enumerate}
\end{proposition}

\begin{proof}
	(i) This is \autoref{lem:cycle-order}.

	(ii) Each \(\beta_i\) has order \(r_i\) by (i). Suppose \(\alpha^M = (1)\). Since the
	\(\beta_i\) commute, \((1) = \alpha^M = \beta_1^M \cdots \beta_t^M\). The \(\beta_i\)
	are disjoint, hence so are their powers, and a product of disjoint permutations can be
	the identity only if every factor is; thus \(\beta_i^M = (1)\) for each \(i\), and then
	\autoref{lem:order-divides} gives
	\(r_i \mid M\) for all \(i\), so \(M\) is a common multiple of \(r_1, \dots, r_t\).
	Conversely, if \(m = \operatorname{lcm}(r_1, \dots, r_t)\), then \(\alpha^m = (1)\).
	Hence \(\alpha\) has order \(m\).

	(iii) Write \(\alpha\) as a product of disjoint cycles and use (ii).
\end{proof}

For example, a permutation in \(S_n\) has order \(2\) if and only if it is either a
transposition or a product of disjoint transpositions. We can now augment
\autoref{tab:cycle-counts}.

\begin{table}[H]
	\centering
	\begin{tabular}{lrrl}
		\toprule
		Cycle structure     & Number & Order & Parity \\
		\midrule
		\((1)\)             & 1      & 1     & Even   \\
		\((1\ 2)\)          & 10     & 2     & Odd    \\
		\((1\ 2\ 3)\)       & 20     & 3     & Even   \\
		\((1\ 2\ 3\ 4)\)    & 30     & 4     & Odd    \\
		\((1\ 2\ 3\ 4\ 5)\) & 24     & 5     & Even   \\
		\((1\ 2)(3\ 4\ 5)\) & 20     & 6     & Odd    \\
		\((1\ 2)(3\ 4)\)    & 15     & 2     & Even   \\
		\midrule
		                    & 120    &       &        \\
		\bottomrule
	\end{tabular}
	\caption{Permutations in \(S_5\).}
	\label{tab:S5}
\end{table}

\subsection{Symmetry}

What do we mean when we say that an isosceles triangle \(\triangle\) is symmetric? Put
\(\triangle = \triangle ABC\) with base \(AB\) on the \(x\)-axis and the \(y\)-axis the
perpendicular bisector of \(AB\). Close your eyes; let \(\triangle\) be reflected in the
\(y\)-axis (interchanging \(A\) and \(B\)); open your eyes. You cannot tell that
\(\triangle\) has been reflected: \(\triangle\) is symmetric about the \(y\)-axis. Had
\(\triangle\) been reflected in the \(x\)-axis it would be obvious. Reflection is a
special kind of \emph{isometry}.

\begin{figure}[H]
	\centering
	\begin{tikzpicture}[scale=1.15, >=stealth]
		\coordinate (A) at (-1.3,0);
		\coordinate (B) at (1.3,0);
		\coordinate (C) at (0,1.8);
		\coordinate (C') at (0,-1.8);
		% reflection of the triangle in the x-axis
		\draw[dashed, thick, BrickRed, fill=BrickRed!6] (A) -- (B) -- (C') -- cycle;
		% the triangle
		\draw[thick, MidnightBlue, fill=MidnightBlue!8] (A) -- (B) -- (C) -- cycle;
		% equal legs
		\draw[MidnightBlue] ($(A)!0.5!(C)!0.1cm!90:(C)$) -- ($(A)!0.5!(C)!0.1cm!-90:(C)$);
		\draw[MidnightBlue] ($(B)!0.5!(C)!0.1cm!90:(C)$) -- ($(B)!0.5!(C)!0.1cm!-90:(C)$);
		\draw[->, gray!80] (2.9,1.0) to[bend left=40]
			node[right=1pt] {\small reflect in \(x\)-axis} (2.9,-1.0);
		\fill (A) circle (1.4pt) node[above left=0pt] {\(A\)};
		\fill (B) circle (1.4pt) node[above right=0pt] {\(B\)};
		\fill (C) circle (1.4pt) node[above right=0pt] {\(C\)};
		\fill[BrickRed] (C') circle (1.4pt) node[below right=0pt] {\(C'\)};
		\node[below left=0pt] at (0,0) {\scriptsize\(O\)};

        % axes
		\draw[->, gray!70] (-2.2,0) -- (2.2,0) node[right] {\(x\)};
		\draw[->, gray!70] (0,-2.4) -- (0,2.5) node[above] {\(y\)};
	\end{tikzpicture}
	\caption{The isosceles triangle \(\triangle ABC\) is symmetric about the \(y\)-axis;
		its reflection in the \(x\)-axis (dashed) is visibly different.}
	\label{fig:isosceles-reflection}
\end{figure}

\begin{definition}[Isometry]
	An \vocab{isometry} of the plane is a function
	\(\varphi\colon \mathbb{R}^2 \to \mathbb{R}^2\) that is \emph{distance preserving}: for
	all \(P = (a,b)\) and \(Q = (c,d)\),
	\[
		\norm{\varphi(P) - \varphi(Q)} = \norm{P - Q},
		\qquad \text{where } \norm{P-Q} = \sqrt{(a-c)^2 + (b-d)^2}.
	\]
\end{definition}

Write \(P \cdot Q = ac + bd\) for the dot product of \(P = (a,b)\) and \(Q = (c,d)\). It
is symmetric and bilinear, and \(P\cdot P = a^2 + b^2 = \norm{P}^2\), where
\(O = (0,0)\) is the origin and \(\norm{P} = \norm{P - O}\). Expanding,
\begin{equation}\label{eq:norm-dot}
	\norm{P-Q}^2
	= (P-Q)\cdot(P-Q)
	= P\cdot P - 2(P\cdot Q) + Q\cdot Q.
\end{equation}

\begin{lemma}[Isometries fixing \(O\) preserve dot products]\label{lem:isometry-dot}
	Let \(\varphi\colon \mathbb{R}^2 \to \mathbb{R}^2\) be a function.
	\begin{enumerate}[(i)]
		\item If \(\varphi\) is an isometry with \(\varphi(O) = O\), then
		      \(\varphi(P)\cdot\varphi(Q) = P\cdot Q\) for all \(P, Q\).
		\item Conversely, if \(\varphi(P)\cdot\varphi(Q) = P\cdot Q\) for all \(P, Q\), then
		      \(\varphi\) is an isometry and \(\varphi(O) = O\).
	\end{enumerate}
\end{lemma}

\begin{proof}
	(i) Putting \(Q = O\) in the definition of an isometry gives
	\(\norm{\varphi(P)} = \norm{\varphi(P) - \varphi(O)} = \norm{P - O} = \norm{P}\), so
	\(\varphi(P)\cdot\varphi(P) = P\cdot P\) for every \(P\). Now expand both sides of
	\(\norm{\varphi(P)-\varphi(Q)}^2 = \norm{P-Q}^2\) using \eqref{eq:norm-dot}:
	\[
		\varphi(P)\cdot\varphi(P) - 2\bigl(\varphi(P)\cdot\varphi(Q)\bigr)
		+ \varphi(Q)\cdot\varphi(Q)
		= P\cdot P - 2(P\cdot Q) + Q\cdot Q.
	\]
	The first terms on the two sides are equal, and so are the last terms; cancelling them
	and dividing by \(-2\) gives \(\varphi(P)\cdot\varphi(Q) = P\cdot Q\).

	(ii) By \eqref{eq:norm-dot} and the hypothesis,
	\begin{align*}
		\norm{\varphi(P)-\varphi(Q)}^2
		 & = \varphi(P)\cdot\varphi(P) - 2\bigl(\varphi(P)\cdot\varphi(Q)\bigr)
		+ \varphi(Q)\cdot\varphi(Q)                                            \\
		 & = P\cdot P - 2(P\cdot Q) + Q\cdot Q = \norm{P-Q}^2,
	\end{align*}
	and taking nonnegative square roots shows that \(\varphi\) is an isometry. Moreover
	\(\norm{\varphi(O)}^2 = \varphi(O)\cdot\varphi(O) = O\cdot O = 0\), so
	\(\varphi(O) = O\).
\end{proof}

\begin{remark}
	The hypothesis \(\varphi(O) = O\) in (i) cannot be dropped. The map
	\(\varphi(U) = U + V\) with \(V = (1,0)\) is an isometry, since
	\(\norm{(P+V) - (Q+V)} = \norm{P-Q}\), but
	\(\varphi(O)\cdot\varphi(O) = V\cdot V = 1 \neq 0 = O\cdot O\).
\end{remark}

Dot products measure angles. Let \(P, Q \neq O\), and write
\(P = \norm{P}(\cos\alpha, \sin\alpha)\) and \(Q = \norm{Q}(\cos\beta, \sin\beta)\) in
polar coordinates. Then
\[
	P\cdot Q
	= \norm{P}\norm{Q}(\cos\alpha\cos\beta + \sin\alpha\sin\beta)
	= \norm{P}\norm{Q}\cos(\alpha - \beta).
\]
The angle \(\theta \in [0,\pi]\) between \(P\) and \(Q\) (that is, \(\angle POQ\)) differs
from \(\pm(\alpha - \beta)\) by a multiple of \(2\pi\); since \(\cos\) is even and
\(2\pi\)-periodic, \(\cos\theta = \cos(\alpha-\beta)\). Hence
\begin{equation}\label{eq:dot-angle}
	P\cdot Q = \norm{P}\norm{Q}\cos\theta,
	\qquad\text{that is,}\qquad
	\cos\theta = \frac{P\cdot Q}{\norm{P}\norm{Q}}.
\end{equation}
Since \(\cos\) is injective on \([0,\pi]\), this determines \(\theta\). In particular,
\(P\) and \(Q\) are orthogonal (\(\theta = \pi/2\)) if and only if \(P\cdot Q = 0\).

\begin{corollary}[Isometries preserve angles]\label{cor:isometry-angle}
	\begin{enumerate}[(i)]
		\item If \(\varphi\) is an isometry with \(\varphi(O) = O\) and \(P, Q \neq O\), then
		      the angle between \(\varphi(P)\) and \(\varphi(Q)\) equals the angle between
		      \(P\) and \(Q\).
		\item More generally, if \(\varphi\) is any isometry and \(P, Q, R\) are points with
		      \(Q, R \neq P\), then
		      \(\angle QPR = \angle\,\varphi(Q)\varphi(P)\varphi(R)\).
	\end{enumerate}
	In particular, every isometry preserves perpendicularity.
\end{corollary}

\begin{proof}
	(i) By \autoref{lem:isometry-dot}, \(\varphi(P)\cdot\varphi(Q) = P\cdot Q\), and, as in
	its proof, \(\norm{\varphi(P)} = \norm{P} \neq 0\) and
	\(\norm{\varphi(Q)} = \norm{Q} \neq 0\). If \(\theta\) and \(\theta'\) are the angles
	between \(P, Q\) and between \(\varphi(P), \varphi(Q)\), then \eqref{eq:dot-angle} gives
	\[
		\cos\theta'
		= \frac{\varphi(P)\cdot\varphi(Q)}{\norm{\varphi(P)}\norm{\varphi(Q)}}
		= \frac{P\cdot Q}{\norm{P}\norm{Q}}
		= \cos\theta,
	\]
	and \(\theta' = \theta\) because both lie in \([0,\pi]\).

	(ii) The angle \(\angle QPR\) is the angle between the vectors \(Q - P\) and \(R - P\).
	Since \((Q-P) - (R-P) = Q - R\), applying \eqref{eq:norm-dot} to \(Q-P\) and \(R-P\)
	gives
	\[
		(Q-P)\cdot(R-P)
		= \tfrac12\bigl(\norm{Q-P}^2 + \norm{R-P}^2 - \norm{Q-R}^2\bigr).
	\]
	The right side involves only distances, as do \(\norm{Q-P}\) and \(\norm{R-P}\), so
	none of them change when \(P, Q, R\) are replaced by \(\varphi(P), \varphi(Q),
	\varphi(R)\). By \eqref{eq:dot-angle}, the two angles have the same cosine, hence are
	equal.

	Finally, perpendicularity means meeting at an angle of \(\pi/2\), which is preserved by
	(ii).
\end{proof}

\begin{notation}
	Denote the set of all isometries of the plane by \(\operatorname{Isom}(\mathbb{R}^2)\).
	Its subset of isometries \(\varphi\) with \(\varphi(O) = O\) is called the
	\vocab{orthogonal group} of the plane, denoted \(O_2(\mathbb{R})\). If \(P \neq Q\) are
	points, let \(L[P,Q]\) denote the line they determine and \(PQ\) the line segment with
	endpoints \(P\) and \(Q\).
\end{notation}

\begin{eg}[Rotations, reflections, translations]\label{eg:isometries}
	\begin{enumerate}[(i)]
		\item Given an angle \(\theta\), the \vocab{rotation} \(R_\theta\) about the origin
		      \(O\) is defined by \(R_\theta(O) = O\), and for \(P \neq O\), by rotating the
		      segment \(PO\) through \(\theta\) (counterclockwise if \(\theta > 0\)). One can
		      rotate about any point of the plane.
		\item The \vocab{reflection} \(\rho_\ell\) in a line \(\ell\), called its
		      \vocab{axis}, fixes each point of \(\ell\); if \(P \notin \ell\), then
		      \(\rho_\ell(P) = P'\) where \(\ell\) is the perpendicular bisector of
		      \(PP'\). Pretending \(\ell\) is a mirror, \(P'\) is the mirror image of \(P\).
		      Now \(\rho_\ell \in \operatorname{Isom}(\mathbb{R}^2)\), and if \(\ell\) passes
		      through the origin, then \(\rho_\ell \in O_2(\mathbb{R})\).
		\item Given a point \(V\), \vocab{translation} by \(V\) is
		      \(\tau_V\colon \mathbb{R}^2 \to \mathbb{R}^2\), \(\tau_V(U) = U + V\). A
		      translation fixes the origin if and only if \(V = O\), so the identity is the
		      only translation that is also a rotation.
	\end{enumerate}
\end{eg}

\begin{lemma}\label{lem:isometry-collinear}
	If \(\varphi\) is an isometry of the plane, then distinct points \(P, Q, R\) are
	collinear if and only if \(\varphi(P), \varphi(Q), \varphi(R)\) are collinear.
\end{lemma}

\begin{proof}
	Suppose \(P, Q, R\) are collinear, with \(R\) between \(P\) and \(Q\), so that
	\(\norm{P-Q} = \norm{P-R} + \norm{R-Q}\). If \(\varphi(P), \varphi(Q), \varphi(R)\) were
	not collinear, they would be the vertices of a triangle, and the triangle inequality
	would give
	\[
		\norm{\varphi(P)-\varphi(Q)} < \norm{\varphi(P)-\varphi(R)} + \norm{\varphi(R)-\varphi(Q)},
	\]
	contradicting that \(\varphi\) preserves distance.

    We can prove the other direction with similar method.
\end{proof}

Every isometry \(\varphi\) is injective: if \(P \neq Q\), then
\(\norm{\varphi(P)-\varphi(Q)} = \norm{P-Q} \neq 0\). Surjectivity is less obvious.

\begin{proposition}\label{prop:isometry-linear}
	Every isometry \(\varphi\) of the plane fixing the origin is a linear transformation.
\end{proposition}

\begin{proof}
	Let \(C_d = \{Q \in \mathbb{R}^2 : \norm{Q - O} = d\}\) be the circle of radius
	\(d > 0\) centred at \(O\). If \(P \in C_d\), then
	\(d = \norm{\varphi(P) - \varphi(O)} = \norm{\varphi(P) - O}\), so
	\(\varphi(C_d) \subseteq C_d\).

	Let \(P \neq O\) and \(r \in \mathbb{R}\). If \(\norm{P - O} = p\), then
	\(\norm{rP - O} = \abs{r}p\), so \(rP \in L[O,P] \cap C_{\abs{r}p}\). Since \(\varphi\)
	preserves collinearity (\autoref{lem:isometry-collinear}),
	\(\varphi(L[O,P]\cap C_{\abs{r}p}) \subseteq L[O,\varphi(P)]\cap C_{\abs{r}p}\), and a
	line meets a circle in at most two points; hence
	\(\varphi(rP) = \pm r\varphi(P)\). To eliminate the minus sign, suppose \(r > 0\). The
	origin lies between \(-rP\) and \(P\), so the distance from \(-rP\) to \(P\) is
	\(rp + p\), whereas the distance from \(rP\) to \(P\) is \(\abs{rp - p}\). Since
	\(rp + p \neq \abs{rp-p}\) and \(\varphi\) preserves distance,
	\(\varphi(rP) \neq -r\varphi(P)\). A similar argument works for \(r < 0\). Therefore
	\(\varphi(rP) = r\varphi(P)\).

	It remains to prove \(\varphi(P+Q) = \varphi(P) + \varphi(Q)\). If \(O, P, Q\) are
	collinear, choose \(U\) on \(L[O,P]\) at distance \(1\) from the origin, so
	\(P = pU\), \(Q = qU\) and \(P + Q = (p+q)U\). Since \(\varphi\) preserves scalar
	multiplication,
	\[
		\varphi(P) + \varphi(Q) = p\varphi(U) + q\varphi(U) = (p+q)\varphi(U) = \varphi(P+Q).
	\]
	If \(O, P, Q\) are not collinear, then \(P + Q\) is the point \(S\) such that
	\(O, P, Q, S\) are the vertices of a parallelogram (parallelogram law). Since \(\varphi\)
	preserves distance, \(O = \varphi(O), \varphi(P), \varphi(Q), \varphi(S)\) are also the
	vertices of a parallelogram, so \(\varphi(S) = \varphi(P) + \varphi(Q)\); that is,
	\(\varphi(P+Q) = \varphi(P) + \varphi(Q)\).
\end{proof}

\begin{corollary}\label{cor:isometry-bijection}
	Every isometry \(\varphi\colon\mathbb{R}^2\to\mathbb{R}^2\) is a bijection, and every
	isometry fixing \(O\) is a nonsingular linear transformation.
\end{corollary}

\begin{proof}
	First assume \(\varphi(O) = O\), so \(\varphi\) is linear by
	\autoref{prop:isometry-linear}. Since \(\varphi\) is injective,
	\(P = \varphi(e_1)\), \(Q = \varphi(e_2)\) is a basis of \(\mathbb{R}^2\), where
	\(e_1 = (1,0)\), \(e_2 = (0,1)\). Hence
	\(\psi\colon aP + bQ \mapsto ae_1 + be_2\) is a well-defined function inverse to
	\(\varphi\); so \(\varphi\) is a bijection and is nonsingular.

	Now let \(\varphi\) be any isometry, say \(\varphi(O) = U\). Then
	\(\tau_{-U}\circ\varphi\) fixes \(O\), so \(\tau_{-U}\circ\varphi = \theta\) is a
	nonsingular linear transformation, and \(\varphi = \tau_U \circ \theta\) is a
	composite of bijections.
\end{proof}

\begin{proposition}\label{prop:isom-group}
	Both \(\operatorname{Isom}(\mathbb{R}^2)\) and \(O_2(\mathbb{R})\) are groups under
	composition.
\end{proposition}

\begin{proof}
	Clearly \(1_{\mathbb{R}^2}\) is an isometry. If \(\varphi'\) and \(\varphi\) are
	isometries, then
	\[
		\norm{(\varphi'\varphi)(P) - (\varphi'\varphi)(Q)}
		= \norm{\varphi(P) - \varphi(Q)}
		= \norm{P - Q},
	\]
	so composition is an operation on \(\operatorname{Isom}(\mathbb{R}^2)\). Each
	\(\varphi\) is a bijection by \autoref{cor:isometry-bijection}, and
	\(\varphi^{-1}\) is also an isometry:
	\(\norm{P-Q} = \norm{\varphi(\varphi^{-1}(P)) - \varphi(\varphi^{-1}(Q))} = \norm{\varphi^{-1}(P) - \varphi^{-1}(Q)}\).
	Composition of functions is always associative. The same proof applies to
	\(O_2(\mathbb{R})\).
\end{proof}

\begin{corollary}\label{cor:isometry-determined}
	\begin{enumerate}[(i)]
		\item If \(O, P, Q\) are noncollinear points and \(\varphi, \psi\) are isometries of
		      the plane with \(\varphi(O) = \psi(O) = O\), \(\varphi(P) = \psi(P)\) and
		      \(\varphi(Q) = \psi(Q)\), then \(\varphi = \psi\).
		\item More generally, if \(A, B, C\) are noncollinear points and \(\varphi, \psi\)
		      are isometries of the plane agreeing on \(A\), \(B\) and \(C\), then
		      \(\varphi = \psi\).
	\end{enumerate}
\end{corollary}

\begin{proof}
	(i) Since \(\varphi\) and \(\psi\) fix \(O\), both are linear transformations by
	\autoref{prop:isometry-linear}. Since \(O, P, Q\) are noncollinear, the list \(P, Q\)
	is linearly independent in \(\mathbb{R}^2\): a nontrivial relation \(aP + bQ = O\)
	would put \(P\) and \(Q\) on a common line through \(O\). As
	\(\dim(\mathbb{R}^2) = 2\), the list \(P, Q\) is a basis, and two linear
	transformations agreeing on a basis are equal.

	(ii) Let \(U = \varphi(A) = \psi(A)\), and set
	\(\varphi' = \tau_{-U}\circ\varphi\circ\tau_A\) and
	\(\psi' = \tau_{-U}\circ\psi\circ\tau_A\). These are composites of isometries, hence
	isometries, and \(\varphi'(O) = \varphi(A) - U = O\); likewise \(\psi'(O) = O\).
	Moreover
	\[
		\varphi'(B - A) = \varphi(B) - U = \psi(B) - U = \psi'(B - A),
	\]
	and similarly \(\varphi'(C - A) = \psi'(C - A)\). The points \(O, B - A, C - A\) are the
	translates of \(A, B, C\) by \(-A\), so they are noncollinear. By (i),
	\(\varphi' = \psi'\), and therefore
	\(\varphi = \tau_U\circ\varphi'\circ\tau_{-A} = \tau_U\circ\psi'\circ\tau_{-A} = \psi\).
\end{proof}

\begin{remark}
	In (i) the hypothesis \(\varphi(O) = \psi(O) = O\) cannot be dropped, and in (ii) two
	points are not enough. If \(O, P, Q\) are noncollinear, the identity and the reflection
	\(\rho_{L[P,Q]}\) both fix \(P\) and \(Q\), yet they are different: the reflection moves
	\(O \notin L[P,Q]\).
\end{remark}

\begin{eg}[Triangles]\label{eg:triangles}
	If \(\triangle\) is a triangle with vertices \(P, Q, U\) and \(\varphi\) is an isometry,
	then \(\varphi(\triangle)\) is the triangle with vertices
	\(\varphi(P), \varphi(Q), \varphi(U)\). If moreover \(\varphi(\triangle) = \triangle\),
	then \(\varphi\) permutes the vertices. Assume the centre of \(\triangle\) is \(O\).
	\begin{itemize}
		\item If \(\triangle\) is \emph{isosceles} with equal sides \(PQ\) and \(PU\), and
		      \(\rho_\ell\) is the reflection with axis \(\ell = L[O,P]\), then
		      \(\rho_\ell(\triangle) = \triangle\); we may describe \(\rho_\ell\) by the
		      transposition \((Q\ U)\). On the other hand
		      \(\rho_{\ell'}(\triangle) \neq \triangle\) for \(\ell' = L[O,Q]\).
		\item If \(\triangle\) is \emph{equilateral}, then also
		      \(\rho_{L[O,Q]}\) and \(\rho_{L[O,U]}\) carry \(\triangle\) to itself
		      (described by \((P\ U)\) and \((P\ Q)\)), and so do the rotations about \(O\)
		      by \(120^\circ\) and \(240^\circ\) (described by \((P\ Q\ U)\) and
		      \((P\ U\ Q)\)).
		\item If \(\triangle\) is \emph{scalene}, then \(\varphi(\triangle) = \triangle\)
		      implies \(\varphi = 1\).
	\end{itemize}
	So an equilateral triangle is ``more symmetric'' than an isosceles triangle, which in
	turn is ``more symmetric'' than a scalene one.
\end{eg}

\begin{definition}[Symmetry group]
	The \vocab{symmetry group} \(\Sigma(\Omega)\) of a figure \(\Omega\) in the plane is the
	set of all isometries \(\varphi\) of the plane with \(\varphi(\Omega) = \Omega\). The
	elements of \(\Sigma(\Omega)\) are called \vocab{symmetries} of \(\Omega\).
\end{definition}

It is straightforward to see that \(\Sigma(\Omega)\) is always a group.

\begin{figure}[H]
	\centering
	\begin{tikzpicture}[scale=1.15, >=stealth]
		% equilateral triangle with its three axes
		\begin{scope}
			\coordinate (P) at (90:1.15);
			\coordinate (Q) at (210:1.15);
			\coordinate (U) at (330:1.15);
			\draw[gray!45, dashed] (P) -- ($(Q)!0.5!(U)$);
			\draw[gray!45, dashed] (Q) -- ($(P)!0.5!(U)$);
			\draw[gray!45, dashed] (U) -- ($(P)!0.5!(Q)$);
			\draw[thick, MidnightBlue] (P) -- (Q) -- (U) -- cycle;
			\fill (P) circle (1.4pt) node[above=2pt] {\(P\)};
			\fill (Q) circle (1.4pt) node[left=2pt]  {\(Q\)};
			\fill (U) circle (1.4pt) node[right=2pt] {\(U\)};
			\fill (0,0) circle (1pt) node[below right=0pt] {\scriptsize\(O\)};
			\node at (0,-2.0) {\(\Sigma(\pi_3) = D_6\), \(\abs{\Sigma(\pi_3)} = 6\)};
		\end{scope}
		% square with its four axes
		\begin{scope}[xshift=4.3cm]
			\coordinate (v0) at (135:1.25);
			\coordinate (v1) at (45:1.25);
			\coordinate (v2) at (-45:1.25);
			\coordinate (v3) at (-135:1.25);
			\draw[gray!45, dashed] (-1.55,0) -- (1.55,0);
			\draw[gray!45, dashed] (0,-1.55) -- (0,1.55);
			\draw[gray!45, dashed] (135:1.55) -- (-45:1.55);
			\draw[gray!45, dashed] (45:1.55) -- (-135:1.55);
			\draw[thick, BrickRed] (v0) -- (v1) -- (v2) -- (v3) -- cycle;
			\fill (v0) circle (1.4pt) node[above left=0pt]  {\(v_0\)};
			\fill (v1) circle (1.4pt) node[above right=0pt] {\(v_1\)};
			\fill (v2) circle (1.4pt) node[below right=0pt] {\(v_2\)};
			\fill (v3) circle (1.4pt) node[below left=0pt]  {\(v_3\)};
			\fill (0,0) circle (1pt) node[below right=0pt] {\scriptsize\(O\)};
			\node at (0,-2.0) {\(\Sigma(\pi_4) = D_8\), \(\abs{\Sigma(\pi_4)} = 8\)};
		\end{scope}
	\end{tikzpicture}
	\caption{Reflection axes of a regular \(3\)-gon and a regular \(4\)-gon.}
	\label{fig:regular-polygons}
\end{figure}

\begin{eg}[Symmetry groups of regular polygons]\label{eg:dihedral-examples}
	\begin{enumerate}[(i)]
		\item A regular \(3\)-gon \(\pi_3\) is an equilateral triangle, and
		      \(\abs{\Sigma(\pi_3)} = 6\), as we saw in \autoref{eg:triangles}.
		\item Let \(\pi_4\) be a square with vertices \(\{v_0, v_1, v_2, v_3\}\), drawn with
		      centre at the origin and sides parallel to the axes. Every
		      \(\varphi \in \Sigma(\pi_4)\) permutes the vertices, and \(\varphi\) is
		      determined by \(\{\varphi(v_i)\}\), so there are at most \(24 = 4!\)
		      symmetries. Not every permutation arises, however: if \(v_i\) and \(v_j\) are
		      adjacent, then \(\norm{v_i - v_j} = 1\), whereas
		      \(\norm{v_0 - v_2} = \sqrt{2} = \norm{v_1 - v_3}\); since isometries preserve
		      distance, \(\varphi\) must preserve adjacency. There are exactly eight
		      symmetries (\autoref{thm:dihedral}): the identity, the three rotations
		      about \(O\) by \(90^\circ, 180^\circ, 270^\circ\), and four reflections with
		      axes \(L[v_0,v_2]\), \(L[v_1,v_3]\), the \(x\)-axis and the \(y\)-axis. This
		      group is the \vocab{dihedral group with \(8\) elements}, denoted \(D_8\).
		\item The symmetry group of a regular pentagon \(\pi_5\) with vertices
		      \(v_0, \dots, v_4\) and centre \(O\) has \(10\) elements: the rotations about
		      the origin by \((72j)^\circ\) for \(0 \leq j \leq 4\), together with the
		      reflections with axes \(L[O,v_k]\) for \(0 \leq k \leq 4\). It is denoted
		      \(D_{10}\).
	\end{enumerate}
\end{eg}

\begin{definition}[Dihedral group]
	A group \(D_{2n}\) with exactly \(2n\) elements is called a \vocab{dihedral group} if it
	contains an element \(a\) of order \(n\) and an element \(b\) of order \(2\) such that
	\[
		bab = a^{-1}.
	\]
\end{definition}

\begin{note}
	The definition is purely algebraic: it does not depend on geometry. If \(n = 2\), then
	\(D_4\) is abelian; if \(n \geq 3\), then \(D_{2n}\) is not abelian. One can show that
	there is essentially only one dihedral group with \(2n\) elements (any two such are
	isomorphic). The name comes from F.~Klein's \emph{dihedron}, two congruent regular
	polygons of zero thickness pasted together.
\end{note}

\begin{theorem}\label{thm:dihedral}
	The symmetry group \(\Sigma(\pi_n)\) of a regular \(n\)-gon \(\pi_n\) is a dihedral
	group with \(2n\) elements.
\end{theorem}

\begin{proof}
	Let \(\pi_n\) have vertices \(v_0, \dots, v_{n-1}\) and centre \(O\). Define \(a\) to be
	the rotation about \(O\) by \((360/n)^\circ\),
	\[
		a(v_i) = \begin{cases}
			v_{i+1} & \text{if } 0 \leq i < n-1, \\
			v_0     & \text{if } i = n-1,
		\end{cases}
		\qquad\text{and}\qquad
		b(v_i) = \begin{cases}
			v_0     & \text{if } i = 0,          \\
			v_{n-i} & \text{if } 1 \leq i \leq n-1,
		\end{cases}
	\]
	where \(b\) is the reflection with axis \(L[O, v_0]\). Clearly \(a\) has order \(n\) and
	\(b\) has order \(2\). There are \(n\) distinct symmetries \(1, a, a^2, \dots, a^{n-1}\),
	and \(b, ab, a^2b, \dots, a^{n-1}b\) are also all distinct by cancellation. If
	\(a^s = a^r b\) with \(0 \leq r \leq n-1\) and \(s = 0, 1\), then \(a^s(v_i) = a^rb(v_i)\)
	for all \(i\); now \(a^s(v_0) = v_s\) while \(a^rb(v_0) = v_{r-1}\), forcing
	\(s = r-1\); but with \(i = 1\) we get \(a^{r-1}(v_1) = v_{r+1}\) while
	\(a^rb(v_1) = v_{r-1}\), a contradiction. So \(a^s \neq a^r b\) for all \(r, s\), and we
	have exhibited \(2n\) distinct symmetries.

	There are no others. We may assume the centre \(O\) of \(\pi_n\) is the origin, so every
	symmetry \(\varphi\) fixes \(O\) and hence is a linear transformation
	(\autoref{prop:isometry-linear}). The vertices adjacent to \(v_0\), namely
	\(v_1\) and \(v_{n-1}\), are the closest vertices to \(v_0\); that is,
	\(\norm{v_i - v_0} > \norm{v_1 - v_0}\) for \(2 \leq i \leq n-2\). So if
	\(\varphi(v_0) = v_j\), then \(\varphi(v_1) = v_{j+1}\) or \(\varphi(v_1) = v_{j-1}\). In
	the first case \(a^j\) agrees with \(\varphi\) on \(v_0\) and \(v_1\), so
	\autoref{cor:isometry-determined} gives \(\varphi = a^j\); in the second case
	\(\varphi = a^j b\) for the same reason. Therefore \(\abs{\Sigma(\pi_n)} = 2n\).

	Finally, \(bab = a^{-1}\): by \autoref{cor:isometry-determined} it suffices to
	evaluate both sides on \(v_0\) and \(v_1\), and
	\(bab(v_0) = v_{n-1} = a^{-1}(v_0)\) while \(bab(v_1) = v_0 = a^{-1}(v_1)\).
\end{proof}

\begin{note}[Symmetry in calculus]
	Symmetry arises in calculus when describing curves \(f(x,y) = 0\):
	\begin{enumerate}[(i)]
		\item symmetry about the \(x\)-axis: the equation is unaltered when \(y\) is replaced
		      by \(-y\) --- in our language, \(\rho_x\);
		\item symmetry about the \(y\)-axis: \(x\) replaced by \(-x\) --- \(\rho_y\);
		\item symmetry with respect to the origin: both replaced --- \(R_{180}\);
		\item symmetry about the line \(y = x\): \(x\) and \(y\) interchanged ---
		      \(\rho_L\) with \(L\) the \(45^\circ\) line.
	\end{enumerate}
	Thus \(f(x,y) = f(x,-y)\) says the graph \(\Gamma\) of \(f(x,y) = 0\) is symmetric about
	the \(x\)-axis, for \((a,b) \in \Gamma\) implies \((a,-b) \in \Gamma\).
\end{note}

\section{Subgroups and Lagrange's Theorem}

A subgroup of a group \(G\) is a subset which is a group under the same operation as in
\(G\). The following definition makes this precise.

\begin{definition}[Closed subset]
	Let \(*\) be an operation on a set \(G\) and let \(S \subseteq G\). We say \(S\) is
	\vocab{closed} under \(*\) if \(x*y \in S\) for all \(x, y \in S\).
\end{definition}

The operation on \(G\) is a function \(*\colon G\times G \to G\); if \(S \subseteq G\),
then \(S \times S \subseteq G\times G\), and \(S\) being closed means
\(*(S\times S) \subseteq S\). For example, \(\mathbb{Z}\) is a subset of the additive group
\(\mathbb{Q}\) closed under \(+\). However \(\mathbb{Q}^\times\) is closed under
multiplication but not under \(+\): both \(2\) and \(-2\) lie in \(\mathbb{Q}^\times\), yet
\(-2 + 2 = 0 \notin \mathbb{Q}^\times\).

\begin{definition}[Subgroup]\label{def:subgroup}
	A subset \(H\) of a group \(G\) is a \vocab{subgroup} if
	\begin{enumerate}[(i)]
		\item \(1 \in H\);
		\item if \(x, y \in H\), then \(xy \in H\) (that is, \(H\) is closed under \(*\));
		\item if \(x \in H\), then \(x^{-1} \in H\).
	\end{enumerate}
\end{definition}

\begin{notation}
	We write \(H \leq G\) to denote that \(H\) is a subgroup of \(G\). Both \(\{1\}\) and
	\(G\) are always subgroups of \(G\). A subgroup \(H\) is \vocab{proper} if \(H \neq G\),
	written \(H < G\), and \vocab{nontrivial} if \(H \neq \{1\}\).
\end{notation}

\begin{proposition}\label{prop:subgroup-is-group}
	Every subgroup \(H \leq G\) of a group \(G\) is itself a group.
\end{proposition}

\begin{proof}
	Axiom (ii) shows \(H\) is closed under the operation of \(G\), so \(H\) has an operation
	(the restriction of \(*\) to \(H\times H\)). This operation is associative, since
	\((xy)z = x(yz)\) holds for all \(x,y,z \in G\), in particular for all
	\(x,y,z \in H\). Axiom (i) gives the identity and axiom (iii) gives inverses.
\end{proof}

\begin{intuition}
	It is \emph{quicker} to check that a subset is a subgroup than to verify the group axioms
	from scratch, because associativity is inherited and need not be verified again.
\end{intuition}

\begin{eg}[Subgroups]\label{eg:subgroups}
	\begin{enumerate}[(i)]
		\item \(O_2(\mathbb{R})\), the isometries fixing the origin, is a subgroup of
		      \(\operatorname{Isom}(\mathbb{R}^2)\). If \(\Omega \subseteq \mathbb{R}^2\),
		      then \(\Sigma(\Omega)\) is also a subgroup of
		      \(\operatorname{Isom}(\mathbb{R}^2)\); and if the centre of gravity of
		      \(\Omega\) is at the origin, then \(\Sigma(\Omega) \leq O_2(\mathbb{R})\).
		\item The four permutations
		      \[
			      \mathbf{V} = \bigl\{(1),\ (1\ 2)(3\ 4),\ (1\ 3)(2\ 4),\ (1\ 4)(2\ 3)\bigr\}
		      \]
		      form a group, because \(\mathbf{V}\) is a subgroup of \(S_4\): \((1) \in
			      \mathbf{V}\); \(\alpha^2 = (1)\) for each \(\alpha \in \mathbf{V}\), so
		      \(\alpha^{-1} = \alpha \in \mathbf{V}\); and the product of any two distinct
		      elements of \(\mathbf{V} - \{(1)\}\) is the third one. One calls
		      \(\mathbf{V}\) the \vocab{four-group}, or \vocab{Klein group}
		      (\(\mathbf{V}\) abbreviates \emph{Vierergruppe}).

		      Note how much work this saves: verifying \(a(bc) = (ab)c\) directly would
		      involve \(4^3 = 64\) equations.
		\item If \(\mathbb{R}^2\) is the plane as an additive abelian group, then any line
		      \(L\) through the origin is a subgroup: choose a point \((a,b)\) on \(L\) and
		      note that \(L\) consists of all scalar multiples \((ra, rb)\).
	\end{enumerate}
\end{eg}

One can shorten the list of items needed to verify that a subset is a subgroup.

\begin{proposition}\label{prop:subgroup-criterion}
	A subset \(H\) of a group \(G\) is a subgroup if and only if \(H\) is nonempty and,
	whenever \(x, y \in H\), then \(xy^{-1} \in H\).
\end{proposition}

\begin{proof}
	If \(H\) is a subgroup, then \(H \neq \varnothing\) since \(1 \in H\); and if
	\(x, y \in H\), then \(y^{-1} \in H\) by (iii), so \(xy^{-1} \in H\) by (ii).

	Conversely, assume \(H\) satisfies the new condition. Since \(H\) is nonempty it
	contains some \(h\); taking \(x = h = y\) gives \(1 = hh^{-1} \in H\), so (i) holds. If
	\(y \in H\), set \(x = 1\) (legitimate now that \(1 \in H\)) to get
	\(y^{-1} = 1y^{-1} \in H\), so (iii) holds. Finally
	\((y^{-1})^{-1} = y\) by \autoref{lem:group-cancellation}, so if \(x, y \in H\), then
	\(y^{-1} \in H\) and \(xy = x(y^{-1})^{-1} \in H\), giving (ii).
\end{proof}

\begin{remark}
	Since every subgroup contains \(1\), one may replace ``\(H\) is nonempty'' by
	``\(1 \in H\)''. If the operation in \(G\) is addition, the condition reads: \(H\) is a
	nonempty subset with \(x, y \in H \implies x - y \in H\).
\end{remark}

\begin{note}[Historical]
	For Galois a group was just a subset \(H\) of \(S_n\) closed under composition.
	A.~Cayley, in 1854, was the first to define an abstract group, mentioning associativity,
	inverses and identity explicitly.
\end{note}

\begin{proposition}\label{prop:finite-subgroup-criterion}
	A nonempty subset \(H\) of a \emph{finite} group \(G\) is a subgroup if and only if
	\(H\) is closed under the operation of \(G\). In particular, a nonempty subset of
	\(S_n\) is a subgroup if and only if it is closed under composition.
\end{proposition}

\begin{proof}
	Every subgroup is nonempty and closed, by axioms (i) and (ii).

	Conversely, assume \(H\) is a nonempty closed subset, so axiom (ii) holds. Then \(H\)
	contains all powers of its elements: there is some \(a \in H\), and \(a^n \in H\) for all
	\(n \geq 1\). As in \autoref{eg:finite-order}, every element of \(G\) has finite
	order, so there is \(m\) with \(a^m = 1\); hence \(1 \in H\) and axiom (i) holds.
	Finally, if \(h \in H\) and \(h^m = 1\), then \(h^{-1} = h^{m-1} \in H\) (since
	\(hh^{m-1} = 1 = h^{m-1}h\)), so axiom (iii) holds.
\end{proof}

\begin{remark}
	This proposition can be false for infinite \(G\): the subset \(\mathbb{N}\) of the
	additive group \(\mathbb{Z}\) is closed under \(+\) but is not a subgroup.
\end{remark}

\begin{eg}[The alternating group]\label{eg:alternating}
	The subset of \(S_n\) consisting of all the \emph{even} permutations is a subgroup,
	because it is closed under multiplication: even \(\circ\) even \(=\) even
	(\autoref{cor:parity-product}). This subgroup is called the \vocab{alternating
		group} on \(n\) letters and is denoted by \(A_n\).
\end{eg}

\begin{note}
	The name comes from polynomials: if
	\(f(x) = (x-u_1)\cdots(x-u_n)\), then \(D = \prod_{i<j}(u_i - u_j)\) changes sign when
	one permutes the roots, since \(\prod_{i<j}[\alpha(u_i) - \alpha(u_j)] = \pm D\). The
	sign \emph{alternates} as various \(\alpha\) are applied; it does not change for
	\(\alpha\) in the alternating group.
\end{note}

\subsection{Cyclic Groups}

\begin{definition}[Cyclic subgroup, cyclic group]
	If \(G\) is a group and \(a \in G\), write
	\[
		\langle a \rangle = \{a^n : n \in \mathbb{Z}\} = \{\text{all powers of } a\};
	\]
	\(\langle a \rangle\) is called the \vocab{cyclic subgroup} of \(G\) generated by
	\(a\). A group \(G\) is \vocab{cyclic} if \(G = \langle a\rangle\) for some
	\(a \in G\), and then \(a\) is called a \vocab{generator} of \(G\).
\end{definition}

It is easy to see \(\langle a\rangle\) is a subgroup: \(1 = a^0 \in \langle a\rangle\);
\(a^n a^m = a^{n+m} \in \langle a\rangle\); \(a^{-1} \in \langle a\rangle\). By
\autoref{eg:group-list}(vi), for every \(n \geq 1\) the multiplicative group
\(\Gamma_n\) of all \(n\)th roots of unity is cyclic, with the primitive \(n\)th root of
unity \(\zeta = e^{2\pi i/n}\) as a generator. A cyclic group can have several different
generators; for instance \(\langle a\rangle = \langle a^{-1}\rangle\).

\begin{proposition}\label{prop:cyclic-generators}
	If \(G = \langle a\rangle\) is a cyclic group of order \(n\), then \(a^k\) is a
	generator of \(G\) if and only if \(\gcd(k,n) = 1\).
\end{proposition}

\begin{proof}
	If \(a^k\) is a generator, then \(a \in \langle a^k\rangle\), so \(a = a^{ks}\) for some
	\(s\). Hence \(a^{ks-1} = 1\), and \autoref{lem:order-divides} gives
	\(n \mid (ks-1)\); that is, \(ks - 1 = tn\) for some integer \(t\), so
	\(sk - tn = 1\) and \(\gcd(k,n) = 1\).

	Conversely, if \(1 = sk + tn\), then \(a = a^{sk+tn} = a^{sk}\) (because \(a^{tn} = 1\)),
	so \(a \in \langle a^k\rangle\). Hence
	\(G = \langle a\rangle \leq \langle a^k\rangle\), and so \(G = \langle a^k\rangle\).
\end{proof}

\begin{corollary}\label{cor:number-of-generators}
	The number of generators of a cyclic group of order \(n\) is \(\phi(n)\), where
	\(\phi\) is the Euler \(\phi\)-function.
\end{corollary}

\begin{proof}
	Immediate from \autoref{prop:cyclic-generators}.
\end{proof}

\begin{proposition}\label{prop:subgroup-of-cyclic}
	Every subgroup \(S\) of a cyclic group \(G = \langle a\rangle\) is itself cyclic. In
	fact \(a^m\) is a generator of \(S\), where \(m\) is the smallest positive integer with
	\(a^m \in S\).
\end{proposition}

\begin{proof}
	We may assume \(S \neq \{1\}\), the statement being obvious otherwise. Let
	\(I = \{m \in \mathbb{Z} : a^m \in S\}\). Then \(0 \in I\), since \(a^0 = 1 \in S\); if
	\(m, n \in I\), then \(a^m a^{-n} = a^{m-n} \in S\), so \(m - n \in I\); and if
	\(m \in I\), \(i \in \mathbb{Z}\), then \((a^m)^i = a^{im} \in S\), so \(im \in I\).
	Since \(S \neq \{1\}\) there is \(a^q \in S\) with \(a^q \neq 1\), so \(q \in I\) and
	\(I \neq \{0\}\). Hence, if \(k\) is the smallest positive integer in \(I\), then
	\(k \mid m\) for every \(m \in I\).

	We claim \(\langle a^k\rangle = S\). Clearly \(\langle a^k\rangle \leq S\). Conversely,
	take \(s \in S\); then \(s = a^m\) for some \(m \in I\), and \(m = k\ell\) for some
	\(\ell\), so \(s = a^{k\ell} \in \langle a^k \rangle\).
\end{proof}

\begin{proposition}\label{prop:order-equals-size}
	Let \(G\) be a finite group and let \(a \in G\). Then the order of \(a\) is the number
	of elements in \(\langle a\rangle\).
\end{proposition}

\begin{proof}
	Since \(G\) is finite, there is \(k \geq 1\) such that \(1, a, a^2, \dots, a^{k-1}\) are
	\(k\) distinct elements while \(1, a, \dots, a^k\) has a repetition; hence \(a^k = a^i\)
	for some \(0 \leq i < k\). If \(i \geq 1\), then \(a^{k-i} = 1\), contradicting the
	original list having no repetitions. Therefore \(a^k = a^0 = 1\), and \(k\) is the order
	of \(a\), being the smallest such positive integer.

	Put \(H = \{1, a, a^2, \dots, a^{k-1}\}\), so \(\abs{H} = k\); it suffices to show
	\(H = \langle a\rangle\). Clearly \(H \subseteq \langle a\rangle\). Conversely, given
	\(a^i \in \langle a\rangle\), the division algorithm gives \(i = qk + r\) with
	\(0 \leq r < k\), and then
	\(a^i = a^{qk+r} = (a^k)^q a^r = a^r \in H\).
\end{proof}

\begin{definition}[Order of a group]
	If \(G\) is a finite group, then the number of elements in \(G\), denoted \(\abs{G}\),
	is called the \vocab{order} of \(G\).
\end{definition}

\begin{abuse}
	The word ``order'' is used in two senses: the order of an element \(a \in G\), and the
	order \(\abs{G}\) of a group. \autoref{prop:order-equals-size} reconciles them:
	the order of \(a\) equals \(\abs{\langle a\rangle}\).
\end{abuse}

\begin{proposition}\label{prop:cyclic-criterion}
	A group \(G\) of order \(n\) is cyclic if and only if, for each divisor \(d\) of \(n\),
	there is at most one cyclic subgroup of order \(d\).
\end{proposition}

\begin{proof}
	Suppose \(G = \langle a\rangle\) is cyclic of order \(n\). We claim
	\(\langle a^{n/d}\rangle\) has order \(d\). Clearly \((a^{n/d})^d = a^n = 1\), and
	\(d\) is smallest with this property: if \((a^{n/d})^r = 1\), then
	\(n \mid (n/d)r\) by \autoref{lem:order-divides}, so \((n/d)r = ns\) for some
	integer \(s\), giving \(r = ds \geq d\).

	For uniqueness, let \(C \leq G\) have order \(d\). By
	\autoref{prop:subgroup-of-cyclic}, \(C\) is cyclic, say
	\(C = \langle x\rangle\) with \(x = a^m\) of order \(d\). Then
	\(1 = x^d = a^{md}\), so \(n \mid md\) by \autoref{lem:order-divides}, say
	\(md = nk\). Therefore \(x = a^m = (a^{n/d})^k\), so
	\(C = \langle x\rangle \subseteq \langle a^{n/d}\rangle\); since both subgroups have
	order \(d\), \(C = \langle a^{n/d}\rangle\).

	Conversely, define a relation on \(G\) by \(a \equiv b\) if
	\(\langle a\rangle = \langle b\rangle\). This is an equivalence relation, and the class
	\([a]\) consists of all the generators of \(C = \langle a\rangle\); denote it by
	\(\operatorname{gen}(C)\). Then
	\[
		G = \bigcup_{C \text{ cyclic}} \operatorname{gen}(C),
		\qquad\text{so}\qquad
		n = \abs{G} = \sum_{C} \abs{\operatorname{gen}(C)},
	\]
	the sum being over all cyclic subgroups of \(G\). But
	\(\abs{\operatorname{gen}(C)} = \phi(\abs{C})\) by
	\autoref{cor:number-of-generators}. By hypothesis \(G\) has at most one cyclic
	subgroup of any given order, so
	\[
		n = \sum_C \abs{\operatorname{gen}(C)} \leq \sum_{d \mid n} \phi(d) = n.
	\]
	Therefore, for each divisor \(d\) of \(n\), there must be a cyclic subgroup \(C\) of
	order \(d\) contributing \(\phi(d)\) to the sum. In particular there is a cyclic
	subgroup of order \(n\), and so \(G\) is cyclic.
\end{proof}

\subsection{Subgroups Generated by a Subset}

\begin{proposition}\label{prop:intersection-subgroup}
	The intersection \(\bigcap_{i \in I} H_i\) of any family of subgroups of a group \(G\) is
	again a subgroup of \(G\). In particular, if \(H, K \leq G\), then \(H \cap K \leq G\).
\end{proposition}

\begin{proof}
	Let \(D = \bigcap_{i\in I} H_i\). Then \(D \neq \varnothing\) because \(1 \in H_i\) for
	all \(i\). If \(x \in D\), then \(x \in H_i\) for all \(i\); each \(H_i\) is a subgroup,
	so \(x^{-1} \in H_i\) for all \(i\) and \(x^{-1} \in D\). Finally, if \(x, y \in D\),
	then \(x, y\) lie in every \(H_i\), hence so does \(xy\), and \(xy \in D\).
\end{proof}

\begin{corollary}\label{cor:generated-subgroup}
	If \(X\) is a subset of a group \(G\), then there is a subgroup \(\langle X\rangle\) of
	\(G\) containing \(X\) that is smallest in the sense that
	\(\langle X\rangle \leq H\) for every subgroup \(H \leq G\) with \(X \subseteq H\).
\end{corollary}

\begin{proof}
	There do exist subgroups of \(G\) containing \(X\); for example \(G\) itself. Define
	\(\langle X\rangle = \bigcap_{X \subseteq H} H\), the intersection of all subgroups
	\(H\) of \(G\) containing \(X\). By \autoref{prop:intersection-subgroup} this is
	a subgroup, and it contains \(X\) because each \(H\) does. If \(H\) is any subgroup
	containing \(X\), then \(H\) is one of the subgroups being intersected, so
	\(\langle X\rangle \leq H\).
\end{proof}

\begin{note}
	There is no restriction on \(X\); in particular \(X = \varnothing\) is allowed. Since
	\(\varnothing \subseteq H\) for every subgroup \(H\), the subgroup
	\(\langle\varnothing\rangle\) is the intersection of \emph{all} subgroups of \(G\), one
	of which is \(\{1\}\); hence \(\langle\varnothing\rangle = \{1\}\).
\end{note}

\begin{definition}[Subgroup generated by a subset]
	If \(X\) is a subset of a group \(G\), then \(\langle X\rangle\) is called the
	\vocab{subgroup generated by \(X\)}.
\end{definition}

\begin{eg}
	\begin{enumerate}[(i)]
		\item If \(G = \langle a\rangle\) is cyclic with generator \(a\), then \(G\) is
		      generated by \(X = \{a\}\).
		\item The symmetry group \(\Sigma(\pi_n)\) of a regular \(n\)-gon is generated by
		      \(a, b\), where \(a\) is the rotation about the origin by \((360/n)^\circ\)
		      and \(b\) is a reflection (\autoref{thm:dihedral}). These generators
		      satisfy \(a^n = 1\), \(b^2 = 1\) and \(bab = a^{-1}\), so
		      \(\Sigma(\pi_n)\) is a dihedral group \(D_{2n}\).
	\end{enumerate}
\end{eg}

\begin{definition}[Word]
	Let \(X\) be a nonempty subset of a group \(G\). A \vocab{word} on \(X\) is either the
	identity element or an element of \(G\) of the form
	\[
		w = x_1^{e_1}x_2^{e_2}\cdots x_n^{e_n},
	\]
	where \(n \geq 1\), \(x_i \in X\) for all \(i\), and \(e_i = \pm 1\) for all \(i\).
\end{definition}

\begin{proposition}\label{prop:words}
	If \(X\) is a nonempty subset of a group \(G\), then \(\langle X\rangle\) is the set of
	all words on \(X\).
\end{proposition}

\begin{proof}
	Let \(W\) be the set of all words on \(X\); we first show \(W \leq G\). By definition
	\(1 \in W\). If \(w = x_1^{e_1}\cdots x_n^{e_n}\) and
	\(w' = y_1^{f_1}\cdots y_m^{f_m}\) lie in \(W\), then
	\(ww' = x_1^{e_1}\cdots x_n^{e_n}y_1^{f_1}\cdots y_m^{f_m}\) is again a word on \(X\),
	so \(ww' \in W\); and
	\(w^{-1} = x_n^{-e_n}x_{n-1}^{-e_{n-1}}\cdots x_1^{-e_1} \in W\). Thus \(W\) is a
	subgroup containing every element of \(X\), and so \(\langle X\rangle \leq W\).

	For the reverse inequality, let \(S\) be any subgroup containing \(X\). Since \(S\) is a
	subgroup it contains \(x^e\) whenever \(x \in X\) and \(e = \pm 1\), and it contains all
	products of such elements; hence \(W \leq S\). Therefore
	\(W \leq \bigcap S = \langle X\rangle\).
\end{proof}

\begin{proposition}[\(\gcd\) and \(\operatorname{lcm}\), group-theoretically]\label{prop:gcd-lcm}
	Let \(a\) and \(b\) be integers and let \(A = \langle a\rangle\) and
	\(B = \langle b\rangle\) be the cyclic subgroups of \(\mathbb{Z}\) they generate.
	\begin{enumerate}[(i)]
		\item If \(A + B = \{\alpha + \beta : \alpha \in A,\ \beta \in B\}\), then
		      \(A + B = \langle d\rangle\), where \(d = \gcd(a,b)\).
		\item \(A \cap B = \langle m\rangle\), where \(m = \operatorname{lcm}(a,b)\).
	\end{enumerate}
\end{proposition}

\begin{proof}
	(i) It is straightforward to check that \(A + B\) is a subgroup of \(\mathbb{Z}\); in
	fact \(A + B\) is precisely the set of all linear combinations of \(a\) and \(b\). By
	\autoref{prop:subgroup-of-cyclic} the subgroup \(A + B\) is cyclic:
	\(A + B = \langle d\rangle\), where \(d\) can be chosen to be the smallest nonnegative
	number in \(A + B\). Thus \(d\) is a common divisor of \(a\) and \(b\), and it is the
	smallest such; that is, \(d = \gcd(a,b)\).

	(ii) If \(c \in A \cap B\), then \(a \mid c\) and \(b \mid c\), so every element of
	\(A \cap B\) is a common multiple of \(a\) and \(b\); conversely every common multiple
	lies in the intersection. By \autoref{prop:subgroup-of-cyclic},
	\(A \cap B = \langle m\rangle\), where \(m\) may be taken to be the smallest
	nonnegative number in \(A \cap B\); that is, \(m = \operatorname{lcm}(a,b)\).
\end{proof}

\subsection{Cosets and Lagrange's Theorem}

Perhaps the most fundamental fact about subgroups \(H\) of a finite group \(G\) is that
their orders are constrained. Certainly \(\abs{H} \leq \abs{G}\), but it turns out that
\(\abs{H}\) must \emph{divide} \(\abs{G}\). To prove this we introduce cosets.

\begin{definition}[Coset]
	If \(H \leq G\) and \(a \in G\), then the \vocab{coset} \(aH\) is the subset
	\[
		aH = \{ah : h \in H\}.
	\]
\end{definition}

Of course \(a = a1 \in aH\). Cosets are usually \emph{not} subgroups: if \(a \notin H\),
then \(1 \notin aH\), for \(1 = ah\) would give the contradiction \(a = h^{-1} \in H\). In
\(*\) notation we write \(a * H = \{a*h : h \in H\}\); in additive notation,
\(a + H = \{a + h : h \in H\}\).

\begin{note}
	The cosets just defined are often called \emph{left} cosets; there are also \emph{right}
	cosets \(Ha = \{ha : h \in H\}\). We shall work almost exclusively with left cosets.
\end{note}

\begin{eg}[Cosets]\label{eg:cosets}
	\begin{enumerate}[(i)]
		\item Consider \(\mathbb{R}^2\) as an additive abelian group and let \(\ell\) be a
		      line through the origin \(O\); as in \autoref{eg:subgroups}(iii), \(\ell\)
		      is a subgroup. If \(\beta \in \mathbb{R}^2\), then the coset \(\beta + \ell\)
		      is the line \(\ell'\) through \(\beta\) \emph{parallel} to \(\ell\), for if
		      \(r\alpha \in \ell\), then the parallelogram law gives
		      \(\beta + r\alpha \in \ell'\) (see \autoref{fig:coset}).
		\item If \(G = S_3\) and \(H = \langle (1\ 2)\rangle\), there are exactly three
		      cosets of \(H\), each of size \(2\):
		      \begin{align*}
			      H            & = \{(1),\ (1\ 2)\}         = (1\ 2)H,   \\
			      (1\ 3)H      & = \{(1\ 3),\ (1\ 2\ 3)\}   = (1\ 2\ 3)H, \\
			      (2\ 3)H      & = \{(2\ 3),\ (1\ 3\ 2)\}   = (1\ 3\ 2)H.
		      \end{align*}
	\end{enumerate}
\end{eg}

\begin{figure}[H]
	\centering
	\begin{tikzpicture}[>=stealth, scale=1.15]
		\coordinate (O) at (0,0);
		\coordinate (a) at (1.6,0.8);
		\coordinate (b) at (-1.0,1.0);
		\coordinate (s) at ($(a)+(b)$);
		\draw[thick, MidnightBlue] (-2.6,-1.3) -- (2.6,1.3) node[right=2pt] {\(\ell\)};
		\draw[thick, BrickRed] (-2.6,0.2) -- (2.6,2.8) node[right=2pt] {\(\beta + \ell\)};
		\draw[gray!55, dashed] (b) -- (s) -- (a);
		\draw[->, gray!85] (O) -- (a);
		\draw[->, gray!85] (O) -- (b);
		\draw[->, thick, BrickRed!85] (O) -- (s);
		\fill (O) circle (1.4pt) node[below=3pt] {\(O\)};
		\fill (a) circle (1.4pt) node[below right=1pt] {\(r\alpha\)};
		\fill (b) circle (1.4pt) node[left=3pt] {\(\beta\)};
		\fill (s) circle (1.4pt) node[above=3pt] {\(\beta + r\alpha\)};
	\end{tikzpicture}
	\caption{The coset \(\beta + \ell\) is the line through \(\beta\) parallel to \(\ell\).}
	\label{fig:coset}
\end{figure}

\begin{observation}
	In these examples different cosets of a given subgroup do not overlap. That is no
	accident: cosets are equivalence classes.
\end{observation}

\begin{proposition}\label{prop:coset-equivalence}
	If \(H \leq G\), then the relation on \(G\) defined by
	\[
		a \equiv b \quad\text{if}\quad a^{-1}b \in H
	\]
	is an equivalence relation, and the equivalence class of \(a \in G\) is the coset
	\(aH\).
\end{proposition}

\begin{proof}
	\emph{Reflexive:} \(a^{-1}a = 1 \in H\). \emph{Symmetric:} if \(a^{-1}b \in H\), then
	\((a^{-1}b)^{-1} = b^{-1}a \in H\), since subgroups are closed under inverses.
	\emph{Transitive:} if \(a^{-1}b, b^{-1}c \in H\), then
	\((a^{-1}b)(b^{-1}c) = a^{-1}c \in H\), since subgroups are closed under
	multiplication.

	For the classes: if \(x \equiv a\), then \(x^{-1}a \in H\), and
	\autoref{lem:coset-basics} gives \(x \in xH = aH\); thus \([a] \subseteq aH\).
	Conversely, if \(x = ah \in aH\), then
	\(x^{-1}a = (ah)^{-1}a = h^{-1}a^{-1}a = h^{-1} \in H\), so \(x \equiv a\) and
	\(x \in [a]\). Hence \([a] = aH\).
\end{proof}

\begin{lemma}\label{lem:coset-basics}
	Let \(H \leq G\) and let \(a, b \in G\).
	\begin{enumerate}[(i)]
		\item \(aH = bH\) if and only if \(b^{-1}a \in H\). In particular, \(aH = H\) if and
		      only if \(a \in H\).
		\item If \(aH \cap bH \neq \varnothing\), then \(aH = bH\).
		\item \(\abs{aH} = \abs{H}\) for all \(a \in G\).
	\end{enumerate}
\end{lemma}

\begin{proof}
	(i) and (ii) are the general facts about equivalence classes applied to
	\autoref{prop:coset-equivalence}: distinct classes are disjoint, and two
	elements are equivalent if and only if their classes coincide. The second statement of
	(i) follows because \(H = 1H\).

	(iii) The function \(f\colon H \to aH\) given by \(f(h) = ah\) is a bijection, with
	inverse \(ah \mapsto a^{-1}(ah) = h\). Therefore \(H\) and \(aH\) have the same number
	of elements.
\end{proof}

\begin{theorem}[Lagrange's theorem]\label{thm:lagrange}
	If \(H\) is a subgroup of a finite group \(G\), then \(\abs{H}\) is a divisor of
	\(\abs{G}\).
\end{theorem}

\begin{proof}
	Let \(\{a_1H, a_2H, \dots, a_tH\}\) be the family of all the distinct cosets of \(H\) in
	\(G\). Then
	\[
		G = a_1H \cup a_2H \cup \cdots \cup a_tH,
	\]
	because each \(g \in G\) lies in the coset \(gH\), and \(gH = a_iH\) for some \(i\).
	Moreover distinct cosets \(a_iH\) and \(a_jH\) are disjoint, by
	\autoref{lem:coset-basics}(ii). It follows that
	\[
		\abs{G} = \abs{a_1H} + \abs{a_2H} + \cdots + \abs{a_tH}.
	\]
	But \(\abs{a_iH} = \abs{H}\) for all \(i\), by \autoref{lem:coset-basics}(iii), and so
	\(\abs{G} = t\abs{H}\).
\end{proof}

\begin{definition}[Index]
	The \vocab{index} of a subgroup \(H\) in \(G\), denoted \([G:H]\), is the number of
	cosets of \(H\) in \(G\).
\end{definition}

When \(G\) is finite, the index \([G:H]\) is the number \(t\) in the formula
\(\abs{G} = t\abs{H}\) above, so that
\[
	\abs{G} = [G:H]\,\abs{H}.
\]
In particular the index \([G:H]\) is also a divisor of \(\abs{G}\).

\begin{corollary}\label{cor:index-formula}
	If \(H\) is a subgroup of a finite group \(G\), then \([G:H] = \abs{G}/\abs{H}\).
\end{corollary}

\begin{proof}
	Immediate from Lagrange's theorem.
\end{proof}

\begin{eg}
	By \autoref{thm:dihedral}, the symmetry group \(\Sigma(\pi_n)\) of a regular
	\(n\)-gon is a dihedral group of order \(2n\). It contains the cyclic subgroup
	\(\langle a\rangle\) of order \(n\) generated by a rotation \(a\), and
	\([\Sigma(\pi_n) : \langle a\rangle] = 2\). Thus there are two cosets:
	\(\langle a\rangle\) and \(b\langle a\rangle\), where \(b\) is any symmetry outside
	\(\langle a\rangle\).
\end{eg}

\begin{corollary}\label{cor:order-divides-group}
	If \(G\) is a finite group and \(a \in G\), then the order of \(a\) is a divisor of
	\(\abs{G}\).
\end{corollary}

\begin{proof}
	By \autoref{prop:order-equals-size}, the order of \(a\) equals the order of the
	subgroup \(H = \langle a\rangle\); now apply Lagrange's theorem.
\end{proof}

This explains why the orders of the elements of \(S_5\) displayed in \autoref{tab:S5} are
divisors of \(120\).

\begin{corollary}\label{cor:am-equals-1}
	If a finite group \(G\) has order \(m\), then \(a^m = 1\) for all \(a \in G\).
\end{corollary}

\begin{proof}
	By \autoref{cor:order-divides-group}, \(a\) has order \(d\) with \(d \mid m\), say
	\(m = dk\). Thus \(a^m = a^{dk} = (a^d)^k = 1\).
\end{proof}

\begin{corollary}\label{cor:prime-order-cyclic}
	If \(p\) is a prime, then every group \(G\) of order \(p\) is cyclic.
\end{corollary}

\begin{proof}
	Choose \(a \in G\) with \(a \neq 1\) and let \(H = \langle a\rangle\). By Lagrange's
	theorem \(\abs{H}\) divides \(\abs{G} = p\); since \(p\) is prime and \(\abs{H} > 1\),
	it follows that \(\abs{H} = p = \abs{G}\), and so \(H = G\).
\end{proof}

\begin{question}[Is the converse of Lagrange's theorem true?]
	Lagrange's theorem says the order of a subgroup of a finite group \(G\) divides
	\(\abs{G}\). If \(d\) divides \(\abs{G}\), must there exist a subgroup of \(G\) of order
	\(d\)?
\end{question}

\begin{remark}
	The answer is \emph{no}: the alternating group \(A_4\) has order \(12\) but has no
	subgroup of order \(6\).
\end{remark}
